\quellseite{331}
\section{Summarium}

Ist eine Historia, wie Moses in der Moabiter Land gestorben, v. 5., und vom HErrn begraben sei, v. 6., und sein Grab hat niemand können erfahren bis auf diesen heutigen Tag, v. 8. Da sollen wir erstlich lernen, wider der Papisten Irrthum: wo Gott wollte, daß man die verstorbenen Heiligen anbeten, ihren Gräbern, Gebeinen oder Asche Ehre thun sollte, würde er Moses, als seines treuesten und trefflichsten Dieners Grab und Gebeine nicht selbst verborgen haben.

Zum andern ist hie zu merken das große Wunderwerk, daß Moses hundert und zwanzig Jahr alt worden, und doch weder am Gesicht, noch sonst am Leibe etwas habe abgenommen. Zu einem Zeugniß, auf daß wir lernen, wer Gott und sein Wort für Augen hat, daß Gott denselbigen mit Leib und Leben erhalten, für Krankheit behüten, und sonst in andere Wege mehr seinen Segen reichlich widerfahren lassen.

\section{Votum}

Also hat euer Liebe das fünfte und letzte Buch Mose gehört, in welchem sonderlich zu merken ist, wie treulich er uns alle vermahnet, daß wir fromm sein, und GOttes Gebot sollen mit Fleiß halten, so werde GOtt auch mit seinem Segen bei uns sein. Wer es aber nicht thut, sondern GOtt und sein Wort verachtet, und sündlich lebet, der soll darum hie und dort gestraft werden. GOtt verleihe seine Gnade, daß wir durch solche Lehre mit Ernst uns bessern und frömmer werden, Amen.

\part{Das Buch Josua}

\chapter{Das 1. Kapitel}

\section{Vorrede}

Euer Liebe haben gehört in den fünf Büchern Mose, daß Gott erstlich im Anfang Himmel und Erde dem Menschen zu Nutz erschaffen hat, und daß der Mensch hernach durch den Teufel im Paradies verführet, in die Sünde und Ungehorsam gefallen sei. Solche sündliche Natur hat sich für und für hernach in der Welt gereget, je länger je mehr, bis Gott mit der Sündflut alle Menschen um der Sünde willen ersäufet hat, aufgenommen den frommen Noah, sammt allem was in der Arch oder Kasten bei ihm war. Von denselbigen acht Menschen ist hernach die neue Welt kommen; dieselbige hat eben der vorigen Welt Exempel nach gelebet, ist gefallen von Gott, seinem Wort und Willen, auf Abgötterei und in allerlei Sünde; denn Fleisch und Blut kann nicht anders, wo der heilige Geist nicht sonderlich hält und behütet. Da erwählet ihm Gott Abraham den Propheten und verhieß ihm, er wollte sein Gott sein, und seines Samens nach ihm; und wollte ihm nicht allein das Land Canaan zu eigen
\quellseite{332}
gehen, sondern auch den gebenedeiten Samen Christum den HErrn, aus seinem Fleisch lassen geboren werden, durch welchen alle Völker auf Erden sollten gesegnet, das ist, von Gottes Zorn und dem Fluch, so sie durch die Sünde verdienet hätten, ledig werden und zu Gnaden kommen, wo sie anders an ihn glaubten. Auf solche Zusage hänget Gott dem Abraham und seinem Samen ein sonder Zeichen, die Beschneidung an, bei welchem sie erkennen sollten, daß sie Gottes Volk wären, und etwas Sonderes für allen andern Völkern auf Erden.

Als nun des frommen Abrahams Nachkommen Hungers halben aus Canaan entweichen mußten in Egypten, und da eine Zeit lang bauseten, aber hernach von den Egyptiern unbillig und hart bedrängt und beschweret worden, da nahm sich Gott allererst ihrer, als seines Volks an, plaget die Egypter, wie das andere Buch Mose zeuget, den König mit aller seiner Macht ersäufet er im rothen Meer, und erlediget also sein Volk.

Von dem an unterscheidet er sie erst recht von allen andern Völkern, nicht allein durch die Beschneidung, sondern auch durchs Gesetz und Eröffnung seines göttlichen Willens, was sie thun und lassen, wie sie ihm auch äußerlich mit Opfern und Anderm dienen sollten. Daß also dies Volk sein eigen Königreich und Priesterthum, oder Kirchenordnung hätte, von Gott selbst gestiftet und geordnet; daher sind in den fünf Büchern Mose manche schöne Predigt, auch seine Exempel, wie man Gott und seinem Wort folgen soll, und wie es Gott entweder belohnet, oder gestraft habe, wenn man nach oder wider sein Wort gethan hat.

Auf diese fünf Bücher folget nun die ganze Historia, wie die Juden, Abrahams Nachkommen, durch Gottes Hülfe das Land eingenommen, und sich darnach im Lande unter den Richtern und Königen gehalten haben, bis sie endlich um ihrer Sünde willen vom Könige aus Assyrien und vom Könige zu Babel, gar aus dem Lande sind vertrieben und weggeführt worden.

Solche Historien, ob sie schon nun längst vergangen sind und zum alten Testament gehören, sind sie doch in mancherlei Weg nütze, daß man sie unter den Christen lesen und lernen soll; denn da findet man nicht allein viel schöne Lieder und Predigten, die sich bisweilen zutragen vom rechten Gottesdienst und allerlei guten Werken, sondern auch die feinsten Exempel, die wohl und fleißig zu merken sind, aus welchem wir zweierlei lernen sollen. Erstlich Gott fürchten, als der die Sünde nicht hat ungestraft lassen hingehen, und darnach auch Gott trauen, als der die Seinen nicht hat stecken lassen, sondern allwege zu gelegener Zeit seine Hülfe mittheilet. Solches ist an ihm selbst nützlich und gut, sintemal ein Christ sein ganzes Leben über stets in der Uebung des Glaubens und Furcht Gottes leben soll. Zum andern ists auch sonderlich tröstlich, daß man in solchen Historien siehet, wie die feinsten, frömmsten und heiligsten Leute bisweilen gestrauchelt, und Unrecht thun, und darnach mancherlei Weise durch Kreuz und Leiden sind geläutert, und doch endlich aus Gottes Gnade solche Sünden ihnen vergeben sind worden.

Zum letzten soll uns je nichts lustigers und lieblichers sein, denn daß wir in solchen Historien lernen mögen, nicht allein uns selbst kennen, und bessern, sondern auch gleich eine Erquickung haben, wenn wir der jetzigen bösen Welt satt und überdrüssig sind, daß wir in die vorige Welt unter die frommen heiligen Leute sehen, und da einen Trost schöpfen, und an ihrem Exempel lernen geduldig sein, und der Welt Bosheit überwinden mit Geduld und Hoffnung eines bessern Lebens, nach diesem sündhaftigen elenden Leben, Amen.

\section{Josua wird in seinem Beruf gestärket}

a Wer nach Gottes Wort sich richtet, der handelt weislich und glücklich; wer nach seinem Kopf führet, der handelt unweislich und vergeblich.
\quellseite{333}
\section{Summarium}

Dies ist eine Historia, wie der HErr den Krieger Josua beruft, daß er soll nach Mose-Abgang des Volks Israel Hauptmann sein; tröstet und lehret ihn, wie er sich halten soll; denn es gehöret ein großer Muth dazu, wer ein Regent sein und den Sachen allenthalben recht thun soll. Darum befehlet Gott mit Fleiß, daß Josua sich fest an Gottes Wort und Gebot halten, und demselbigen folgen soll. Laß das Gesetz des HErrn (spricht er V. 8.) nicht von deinem Munde kommen, alsdenn wird dir gelingen in allem das du thust, und wirst weislich handeln können. Solches soll insonderheit die Obrigkeit zu Herzen nehmen, so sie anders glücklich und recht regieren will; denn wo man Gottes Wort nicht folgen, sondern dasselbige verfolgen will, da wird weder Glück noch Heil sein können.

\section{Votum}

Der treue, barmherzige GOtt erhalte uns sein heiliges Wort, und regiere durch den heiligen Geist Obrigkeit und Unterthanen, daß sie seinem Wort gehorsamlich folgen, ihm zu großem Lob, Ehr und Preis, und den Menschen zu gemeiner heilsamen Wohlfahrt, Amen.

\chapter{Das 2. Kapitel}

\section{Vorrede}

Gott will ein gnädiger Vater sein allen Armen und Gläubigen, und soll ihm kein Mensch zu geringe noch zu verachtet sein. Wiederum will Gott ein zorniger Richter sein über die Gottlosen und Ungläubigen, und soll ihm keiner zu groß, zu mächtig, zu klug und zu reich sein; beides bezeugen die geistlichen und weltlichen Historien alter Geschichte und die tägliche Erfahrung. Derwegen wollen wir uns demütigen und in wahrem Glauben und Gehorsam an Gott halten, so genesen wir in Gefahr, in Noth und Tod. Wer aber Gott und seinem Wort nicht trauet noch folget, der wird auch den Lohn haben, aber mit Schand und Schaden. Dies soll auch erinnern das 2. Kapitel.

\section{Die Kundschafter werden bei der Rahab errettet}
\quellseite{334}
\section{Summarium}

Sie ist ein schön Exempel eines feinen Glaubens, welchen die Epistel an die Ebräer auch rühmet: daß Rahab für gewiß achtet, daß die Kinder Israel Gottes Volk sind und sollen wider die Heiden siegen, und ihr Land einnehmen, weil Gott in Egypten und am rothen Meer so große Wunderzeichen mit ihnen gethan habe. Darum verbirgt sie die Kundschafter, V. 6., und wird durch sie auch erhalten, da die ganze Stadt dagegen keinen Glauben an Gottes Wort

19 *
\quellseite{335}
hatte, und derhalben alles darin mußte erwürget werden.

\section{Votum}

Gott, der Vater unsers Herrn Jesu Christi, stärke unsern schwachen Glauben, und helfe, daß wir durch den Glauben überwinden den Teufel, die Welt, auch unser eigen Fleisch und Blut, und daß wir alle an Leib und Seel erhalten werden, zeitlich und ewig, Amen.

\chapter{Das 3. Kapitel}

\section{Vorrede}

Gleich wie die levitische Bundeslade zur Zeit Josua mitten im Jordan gestanden, den Strom desselben aufgehalten, und den Israeliten, nach ihrer schweren und gefährlichen vierzigjährigen Pilgrimschaft einen sichern und freien Paß und Gasse ins gelobte Land verschaffet: also da unser Herr Christus hat sollen getauft werden, und in sein Predigtamt treten, kömmt er als der lebendige Gnadenthron, und stellet sich selbst mitten in den Jordan, und will unser Weg und Steg zum ewigen Vaterland sein, und schützen und aufhalten alle Schüsse und Güsse der Flut göttliches Zorns, und des Wassers der Trübsalen abzuwehren, daß uns die Ströme des rothen Drachens und Satans nicht verschw emmen, noch die Meereswellen der Widerwärtigkeit allzu hart über unsere Seele gehen, und uns gar tilgen sollen, so wir nur auf ihn sehen, und uns an ihn in wahrem Glauben halten. Das sollen wir uns auch erinnern bei dem 3. Kapitel.

\section{Israel geht trockenes Fußes durch den Jordan}

a Diese Geschichte bedeutet, daß uns Christus durch die Apostel vorgetragen im Evangelio, leitet ins Himmelreich durch den trockenen Jordan, der doch diesmal am vollesten war. Das ist, daß Gesetz, das uns mit Werken ängstet und treibet, höret auf durch Evangelium, daß unsere Gewissen frei, fröhlich und sicher für ihm sind und allein Christo im Glauben folgen.
\quellseite{336}
\section{Summarium}

Ist eine Historia des herrlichen Wunderwerks, daß Gott den Jordan still stehen und sein Volk mit trockenem Fuß ließ hindurch gehen V. 15., auf daß sie an solchem Wunderwerk lernten, Gott wollte mit ihnen sein und sie derhalben desto getroster hinan gingen, und das Land einnähmen.

Die geistliche Deutung ist diese, daß Christus gleich wie die Archa uns wird fürgetragen im Evangelio. Demselben gehen wir nach im Glauben durch den Jordan ins gelobte Land, das ist, durch dies armselige Leben ins ewige Leben, da sonst, wo Christus nicht wäre, der Jordan, das ist, die Welt sammt dem Teufel, Sünd und allem Unglück uns würde ersäufen.

\section{Votum}

Unserm HErrn JEsu Christo sei Lob, Preis und Dank für alle seine Wohlthaten; der wolle auch heute fürgehen allen, die durch den Jordan des Trübsals wandern müssen, und wolle uns allesamt in diesem Elend und im zeitlichen Tode zum ewigen Leben erhalten, Amen.

\chapter{Das 4. Kapitel}

\section{Vorrede}

Die christliche Kinderzucht ist leider jetziger Zeit gar ungemeine. Weh aber über alle, weh allen Eltern, die ihre Kinder nicht in der Zucht und Vermahnung zu dem HErrn auferziehen. Denn kein Löwe ist so grausam, kein Bär so grimmig, kein Wolf so wüthend, keine Schlange so giftig, keine Gluckhenne so hungrig, sie versorgen ihre Jungen nach ihrer Art, mit aller Nothdurft. Sollte denn der Mensch, der doch die alleredelste, herrlichste und die vernünftigste Creatur nach den Engeln sein soll, sein Fleisch und Blut mit nöthiger Zucht verwahrlosen und verlassen? Derwegen die Eltern, und wem mehr die Kinderzucht befohlen ist, sollen also die Jungen erziehen, daß sie Gott zu Ehre, ihrem Nächsten zu Nutz, und ihnen selber zur Wohlfahrt leben, auf daß sie auch selber an den Kindern Ehre und Freude erleben, und nach diesem betrübten Leben mit ihnen ewige, freudenreiche Gemeinschaft, sammt der hochgelobten Dreifaltigkeit, allen auserwählten Engeln und Christgläubigen haben mögen; und sollen wissen, so manche Mühe sie mit der Kinderzucht haben, so manchen Gottesdienst begehen sie damit, wo ferne sie aus dem Glauben solche Werke thun. Dies lehret auch das 4. Kapitel.

\section{Denckzeichen des wunderbaren Durchgangs durch den Jordan}

a Also das ist, die Priester, so die Lade trugen, stunden, wie sie geordnet waren, und thaten nichts anders.
\quellseite{337}
\section{Summarium}

Ist ein Befehl, wie man zwölf Steine zum Gedächtniß des großen Wunderwerks soll aufrichten, v. 3.; denn Gott will sonderlich, daß man seine Wohlthat nicht vergessen, sondern dieselbigen dem jungen Volk fleißig einpflanzen, und es dabei erhalten soll, auf daß man in andern fürfallenden Nöthen auf seine Hülfe bauen und ihm trauen lerne. Darum sollen die Eltern hie sich ihres Amts erinnern, und mit sonderlichem Fleiß die Kinder zu Gottes Wort und Furcht halten.

\section{Votum}

Helfe der treue GOtt, daß die Eltern in der Kinderzucht fleißig und treu, und daß die Kinder gehorsam seyn und sich gerne ziehen lassen, auf daß die Eltern und Kinder GOtt zeitlich und ewig mit einander loben und preisen, Amen.

\chapter{Das 5. Kapitel}

\section{Vorrede}

Gott verläßt niemand, man verlasse ihn nur nicht; denn je größer Noth, je näher GOtt ist, wenn man ihn in wahrem Glauben und Vertrauen auf seine Verheißungen anruft um Trost und Hülfe in der Noth. Derwegen wollen wir auch mit Freuden fleißig thun, was einem jeden in seinem Beruf gebühret, und wollen beten, daß Gott helfen, schützen und erretten wolle. So werden wir auch Gottes Segen haben, und von ihm behütet und erhalten werden. Denn also tröstet Gottes Wort im fünften Buch Mose im 4. Kapitel V. 31.: Der HErr dein Gott ist ein barmherziger GOtt, er wird dich nicht verlassen, noch verderben, wird auch nicht vergessen des Bundes, den er deinen Vätern geschworen hat. Davon höret auch an das 5. Kapitel.

\section{Beschneidung und Passah im gelobten Lande verrichtet}

a (Steinerne Messer) scharf, als am Stein gewetzt, Ps. 89, 44. Avertisti petram gladii ejus. Die Schärfe seines Schwerts ist stumpf, daß es nicht schneidet.
\quellseite{338}
b) Sangen) am Feuer gedörrte oder versengte Aehren, tostatas spicas.

\section{Summarium}

Dies ist erstlich eine Anzeige, was Gott mit dem Wunderwerk verrichtet habe, daß der Jordan ist ausgetrocknet, nämlich, daß sein Volk dadurch getröstet, und den Feinden das Herz genommen würde. Eben um solcher Ursache willen erscheinet ein Engel Josua insonderheit, V. 13., auf daß er ja allenthalben sehe und merkte, daß Gott wolle bei ihm seyn, und er deßhalben desto freudiger und getroster seinem Beruf soll nachkommen.

\section{Votum}

Gott lasse ihm alle christliche Obrigkeit befohlen sein, verleihe ihnen Hülfe und den Unterthanen Schutz, und segne auch die Werke unserer Hände zu seinem Lob, und unsers Nächsten Wohlfahrt, um Jesu Christi, unsers Herrn, willen, Amen.

\chapter{Das 6. Kapitel}

\section{Vorrede}

Gott widerstrebet den Gottlosen und hilft den Gläubigen. Denn obschon die Kinder Gottes auch ihre Mühseligkeit und Kreuz auf Erden haben, und oft mehr und größeres, denn die Sünden-Knechte; dennoch folgt ihnen ein gewünschter Ausgang, die Gottlosen aber gehen endlich gar zu Grunde und Boden. Davon sagt David im 75. Psalm V. 9. also: Der HErr hat einen Becher in der Hand und mit starkem Wein voll eingeschenkt, und schenket aus demselben; aber die Gottlosen müssen alle trinken und die Hefen aussaufen. Derwegen wer Gott mit seinem Segen auf seiner Seite haben will, der thue Buße, bete und lebe nach seinem Wort und Willen unsträflich. Dies erinnert auch das 6. Kapitel.

\section{Eroberung und Zerstörung der Stadt Jericho}
% Scan page 339: No non-biblical apparatus, gloss, or structural title/synopsis retained in OCR blocks.
\quellseite{340}
\section{Summarium}

Ist eine Historia, daß die Stadt Jericho jenseit des Jordans am ersten genommen, und mit allem verbrannt sei worden, V. 2. Hie soll man erstlich lernen, wie Gott solch greulich Urtheil über diese Stadt um der Sünden willen hat gehen lassen, daß Alt und Jung, Mann und Weib erstochen ward, V. 21., auf daß wir uns für Sünden hüten, und dergleichen Urtheil nicht auch verdienen. Zum andern, daß die Mauern von ihnen selbst umfallen, und das Volk Israel doch nicht mehr thut, denn daß es um die Stadt gehet mit der Arche, V. 30., solches ist ein Exempel durch welches Gott sie hat wollen lehren, wie er für sie streiten und ihnen beistehen wolle. Wie denn aller Sieg von Gott kommt, und alles an Gottes Segen liegt.

\section{Votum}

Der allmächtige, barmherzige GOtt streite heute auch für seine Gemeine, und segne das Kirchen-, weltliche und Hausregiment, helfe doch, daß ein jeder in seinem Beruf christlich und ehrbar lebe, Amen.

\chapter{Das 7. Capitel}

\section{Vorrede}

Gottes Wort lehret es, die Historien und aller Welt Geschichte bezeugen, und die tägliche Erfahrung bestätiget, daß Gott die Sünder endlich ernstlich und heftig strafet. Aber, ach der Welt Blindheit und Unsinnigkeit, die sich dies nicht überreden und bessern will lassen. Derwegen wollen wir unser und der Unsern verschonen, wollen die Sünde hassen und lassen, und Tugend lieben und üben, daß solches sei Gott ein Wohlgefallen, den heiligen Engeln eine Freude, unserm Nächsten Nutz, uns zur Wohlfahrt diene, und daß den Teufeln eine Pein und Marter sei ewiglich, Amen. Darauf höret ferner an das 7. Capitel.

\section{Achan wegen Diebstahls bestraft}
\quellseite{341}
\section{Summarium}

Hie ist ein Exempel, wie Gott bisweilen an einem ganzen Haufen strafet, was etliche insonderheit haben unrecht gethan, V. 35., auf daß die weltliche Obrigkeit lerne ihres Amts fleißig warten, und das Uebel mit höchstem Ernst strafen, denn thut mans nicht und läßt die Strafe unterwegen, so muß hernach der ganze Haufe darum herhalten, wie man siehet.

a Spangen, wie eine Zunge gestaltet.

\section{Summarium}

Ist ein sonderlich Exempel, in welchem wir sehen, wie hart und mit großem Ernst Gott die Sünde strafet, und endlich an den Tag will bringen, man berge und decke sie so lang man immer wolle. Denn weil Gott verboten hatte, daß niemand vom Raub etwas sollte nehmen, und Achan dennoch Gold und Silber genommen und vergraben hatte, mußte es an Tag kommen, und er mit Weib und Kindern, mit Vieh und aller seiner Habe verbrennet werden, V. 26. Solche heftige Strafe soll uns allen eine Warnung sein, daß wir uns für Sünden hüten und wider Gottes Wort nichts thun noch fürnehmen sollen, auf daß uns dergleichen nicht auch widerfahre.

\section{Votum}

Regiere der treue GOtt die Obrigkeit, daß sie die Sünden ernstlich strafe, und uns regiere GOtt, daß wir uns für Sünden fleißig hüten, auf daß wir allem Unglück entfliehen, und bestehen mögen am jüngsten Tage für des Menschen Sohn, Amen.

\chapter{Das 8. Kapitel}

\section{Vorrede}

Ein jeder Beruf und Handthierung hat seine Mittel, dadurch sie sollen getrieben werden. Diese soll ein jeder fleißig üben und gebrauchen, und den Segen dazu allein von Gott bitten und hoffen, nach dem gemeinen Sprüchwort: Bete und arbeite. Denn wer sein Vertrauen auf die Mittel allein setzt, als auf seinen Fleiß und Arbeit, der hat wohl viel Mühe, aber seinen Nutzen hat er davon, wie Gottes Wort und die Erfahrung bezeuget. Dagegen der Segen des HErrn
\quellseite{342}
macht reich ohne Mühe, sagt Salomon Sprüche 10, 22. Und St. Paulus spricht von denen, die mit Gott und Ehren ihres Berufs abwarten: Eure Arbeit ist nicht vergeblich in dem HErrn, 1 Cor. 15, 58. Davon unterrichtet auch das 8. Kapitel.

\section{Belagerung und Eroberung der Stadt Ai}
\quellseite{343}
\section{Summarium}

Ist eine Historia, wie Ai und Bethel von Josua gewonnen sind, V. 2. Daß aber der HErr befiehlt, daß Josua mit List sie aus der Stadt locken, und darnach von hinten zufallen sollen, ist eine Anzeigung, daß, ob wohl der Sieg von Gott kömmt, und alles am Segen Gottes liegt, dennoch wir allerlei Mittel, so von Gott verordnet sind, dahin brauchen sollen, und gleichwohl unsere Hoffnung auf Gott und nicht auf die Mittel setzen. Also soll ein Kriegsmann seinen Vortheil nicht begeben, Büchsen und Harnisch brauchen, aber auf solche Mittel sich nicht verlassen, sondern seine Zuversicht allein zu Gott haben.

\section{Votum}

Gott regiere uns alle, daß wir ihm allein die Ehre geben, in allem unsern Vornehmen und Thun, und segne und behüte einen jeden in seinem christlichen Beruf und Stande, daß es zu seinem Lob und gemeinem Nutz zum Besten diene und gelange, Amen.

\chapter{Das 9. Kapitel}

\section{Vorrede}

Die Alten nennen Eidschwur eine Arznei; denn wie man die Arznei nicht oft und ohne sonderer Noth nicht gebrauchen soll, also sollen wir ohne große erhebliche Noth des Eides nicht gebrauchen. Wenn er aber zur Stärkung der Wahrheit, Gott zu Ehren, uns und dem Nächsten zur Wohlfahrt, auch zur Erhaltung Friedens und Einigkeit und Rettung der Unschuld dienet, und ist nicht wider Gottes Wort, so kann und soll man die Eide gebrauchen und thun, und sonst soll man sich für unnützen und unnöthigen Schwüren hüten. Denn schwören ohne Noth (spricht man) ist der Seelen Tod. Derwegen warnet uns Sirach ernstlich dafür und sagt im 23. Kapitel V. 9.: Gewöhne deinen Mund nicht zum Schwören, denn gleich als ein Knecht, der oft gepräget wird, nicht ohne Striemen ist, also kann einer der oft schwöret, nicht ohne Sünden sein. Dies lehret auch das 9. Kapitel.

\section{Die Gibeoniter kommen mit List in den Bund}
% Scan page 344: No non-biblical apparatus, gloss, or structural title/synopsis retained in OCR blocks.
\quellseite{345}
\section{Summarium}

Hie ist eine sonderliche Historia, in welcher wir sehen, wie man den geschworenen Eid soll halten, obs gleich mit unserm Nachtheil geschehen muß, und ob wir übervortheilt sind. Das also keine Ausflucht ist, was man geschworen und zugesagt hat, das soll man halten, ohne so es wider Gott und sein Wort ist.

Darneben siehet man auch hie, wie die, so durch List mit ihrem Vortheil den Eid suchen, ungestraft nicht bleiben. Denn die Gibeoniter sammt allen ihren Nachkommen müssen diese Strafe tragen, daß sie Holzhauer und Wasserträger unter den Kindern Israel müssen bleiben, V. 27. Hie mag man bedenken, was für eine Strafe werde erfolgen, wo man einen falschen Eid schwört, oder den geschworenen Eid um Geldes oder Guts willen bricht. Denn weil der Name Gottes zum Eid gebraucht, und der Eid dadurch für kräftig gehalten wird, will Gott seinen Namen von niemand unehren, noch die Leute dadurch betrügen lassen.

\section{Votum}

Der barmherzige GOtt regiere uns, daß wir im Thun und Lassen nach seinem Wort uns halten, und also nach dem Eide, den wir ihm in der Taufe gethan, christlich leben, damit es zu seinen Ehren und zu unserer Wohlfahrt gelange, Amen.

\chapter{Das 10. Kapitel}

\section{Vorrede}

Die Gott kennen, die ihn fürchten, zeitlich und ewiglich wohl zu thun verheißt, also dräuet er wiederum zeitlich und ewig zu strafen, die so ihn verachten. Von Gottesfürchtigen sagt David also Psalm 33, 18.: Siehe des HErrn Auge stehet auf die, so ihn fürchten, die auf seine Güte hoffen, daß er ihre Seele errette vom Tode, und ernähre sie in der Theurung. Und abermals spricht er Psalm 147, 11.: Der HErr hat Gefallen an denen, die ihn fürchten, die auf seine Güte hoffen. Von Gottesverächtern sagt Gott Mal. 3, 5. also: Ich will zu euch kommen und euch strafen, und will ein schneller Zeuge sein wider die, so mich nicht fürchten. Und im Buch der Weisheit Cap. 1, 8. stehet diese Dräuung: Der kann nicht verborgen bleiben, der das Unrecht redet, und das Recht, so ihn strafen soll, wird sein nicht fehlen; denn des Gottlosen Anschläge müssen vor Gericht, und seine Reden müssen für den HErrn kommen, daß seine Untugend gestrafet werde. Davon höret auch an das 10. Kapitel.

\section{Wunderbarer Sieg Josuas wider die Amoriter}
\quellseite{346}
\section{Summarium}

Ist eine sonderliche Historia und ein herrlich großes Wunderwerk, daß Gott der Stimme Josua, des einigen Mannes gehorchet und läßt Sonne und Mond einen ganzen Tag am Himmel stehen, daß sie nicht fortgehen, V. 13.; daß also das Volk desto mehr Zeit hätte, sich an seinen Feinden zu rächen. Solches Wunderwerk dienet dahin, daß man lerne, wie Gott über den Seinen, die ihn fürchten, halten wolle, und thun was sie begehren, wie der 145. Psalm V. 19. saget. Neben dem, daß auch solch Wunderwerk dem Volk einen Muth, und den Feinden einen Schrecken macht, weil sie sahen, daß Gott mit Josua war.
\quellseite{347}
\section{Summarium}

Ist eine Historia, wie Josua etliche Könige und Städte geschlagen und vertilget habe. Da sehet man, wo Gott mit seinem Segen ist, daß alles schleunig von der Hand gehet. Weil aber Gott seinen Segen allein den Frommen hat zugesagt, sollen wir uns seines Worts und Willens mit Ernst halten, so wird es uns auch glücklich gehen. Wie denn hie allenthalben von Josua gerühmet wird, er hab alles verbannet und ausgerottet, wie der HErr ihm geboten hat. Darum ist auch Glück bei ihm.

\section{Votum}

GOtt, unser lieber, treuer Vater, regiere und leite uns, daß wir seinen heiligen Geboten folgen, und segne und behüte uns an Leib und Seele, in allen Ständen, zu zeitlicher und ewiger Wohlfahrt, um JEsu Christi willen, Amen.

\chapter{Das 11. Capitel}

\section{Vorrede}

Gott pfleget den Verzug der Strafe mit Schwerheit der Strafe einzubringen, das ist, je länger Gott verzucht die Gottlosen zu strafen, je greulicher er sie darnach, wenn sie ihr Maß gefüllet, und er seine Zeit ersehen, zu strafen pfleget. Dieses sollte uns genug Ursache sein, christlich zu leben und für Sünden sich zu hüten. Aber die Welt will gestraft sein, darum ringet sie nach Schlägen und mißbraucht Gottes Langmuth und Geduld erschrecklich zu greulichen Sünden. Derwegen wird sie gestraft, und wird mehr und größere Strafe leiden müssen. Gott erbarme sich unser gnädiglich, Amen. Darauf höret auch davon an das 11. Kapitel.

\section{Ausrottung der Canaaniter}
\quellseite{348}
\section{Summarium}

Das ist ein fein Exempel, in welchem man siehet, wie Gott den Gottlosen eine Zeit lang zusieht, aber wenn sie sich nicht wollen bessern, und das Stündlein kommt, müssen sie ihnen selbst Ursache geben zu ihrem Verderben. Wie Josua hie sagt: Es geschah aber vom HErrn, daß ihr Herz verstocket wäre, daß sie sich unterstünden mit Israel zu streiten, V. 26., das ist, ihre Sünden hattens verdient und waren reif zur Strafe, darum mußten sie selbst Ursache dazu geben. Also stehet man, daß kein größerer Vortheil im Kriege wider die Feinde ist, denn so sie durch ihre Sünde Gott wider sich haben; wie heutiges Tages des Türken höchste Macht und großer Vortheil ist, daß die Obrigkeit zum Theil Gott und sein Wort verfolget, und sonst auch kräftlich ist, und wir insgemein so ruchlos sind in Geiz, Unzucht, Büberei, Haß, Neid und anderen Sünden, und sehr wenig sich an Gottes Wort bessern, da muß denn wohl Unglück folgen.

\section{Votum}

Der treue GOtt erleuchte unsere Augen und Herzen, daß wir uns an anderer Leute Unglück spiegeln, von Sünden abstehen, gerecht und heilig leben, auf daß wir nicht auch zeitlich und ewig gestraft werden, Amen.
\quellseite{349}
\chapter{Das 12. Kapitel}

\section{Besiegte Könige}

\section{Summarium}

Josua wiederholet hie in einer Summa alle Könige und Städte, so Mose und Josua geschlagen und eingenommen haben, nämlich, ein und dreißig Könige, V. 24. An solcher Summa mag man lernen, wie die Sünde und das gottlose Wesen einen Riß in der Welt und großen Regimenten machet, wie es Krieg und alles Unglück erreget. Darum sollen wir in Gottesfurcht wandeln, und seines Willens und Wortes uns fleißigen, so will er uns für solchem Unglück behüten.

\chapter{Das 13. Kapitel}

\section{Austheilung des Landes wird vorgenommen}
% Scan page 350: No non-biblical apparatus, gloss, or structural title/synopsis retained in OCR blocks.
\quellseite{351}
\chapter{Das 14. Capitel}

\section{Von Calebs Erbtheil}

a Kirjath Arba) das ist die Arbastadt, von Arba einem Riesen also genannt, wie wir Deutschen sagen, Karlstadt, Arnstadt etc.

\section{Summarium}

Ist ein fein Exempel wie Caleb, Josua Gesell, der Kundschafter einer, darum, daß er dem Gott Israel treulich gefolget hatte, von Gott dermaßen ist behütet worden, daß er sich rühmet, er sei fünf und achtzig Jahr alt, und sei dennoch allerdings so stark, als da er vierzig alt sei gewesen, V. 11. Wer aber ein gottlos Leben führet, und nicht wandelt in der Furcht GOttes, da träuet die Schrift, daß dieselbigen ihre Tage nicht auf die Hälfte bringen werden; wie wir zu unserer Zeit Exempel genug, und schier ein jeder in seinem eigenen Hause und an ihm selbst sehen.

\section{Votum}

GOtt, der Vater aller Barmherzigkeit, behüte uns für gottlosem Thun und Wandel, und helfe, daß unser Leben ihm ein wohlgefällig Leben sei, und daß wir Segen von ihm haben an Leib und Seele, in diesem und künftigem Leben, Amen.

\chapter{Das 15. Capitel}

\section{Grenze und Städte des Stammes Juda}
% Scan page 352: No non-biblical apparatus, gloss, or structural title/synopsis retained in OCR blocks.
\quellseite{353}
a Denn David hat sie hernach mit der Zeit vertrieben, 2 Sam. 5, 7.

\chapter{Das 16. Kapitel}

\section{Grenzstätte der Kinder Ephraims}

\chapter{Das 17. Kapitel}

\section{Vom Erbtheil der andern Hälfte des Stammes Manasse}
\quellseite{354}
\chapter{Das 18. Capitel}

\section{Grenzstätte des Stammes Benjamin}
\quellseite{355}
\chapter{Das 19. Kapitel}

\section{Der übrigen sechs Stämme und Josua's Erbtheil}
% Scan page 356: No non-biblical apparatus, gloss, or structural title/synopsis retained in OCR blocks.
\quellseite{357}
\chapter{Das 20. Kapitel}

\section{Verordnung der Freistädte}

\section{Summarium}

Ist eine Ordnung mit den Freistädten, V. 2., auf daß wer unwissend oder wider seinen Willen hätte einen Todtschlag gethan, nicht übereilet würde, wie zuvor im Mose weiter davon gesagt ist.

\chapter{Das 21. Kapitel}

\section{Versorgung der Leviten}
% Scan page 358: No non-biblical apparatus, gloss, or structural title/synopsis retained in OCR blocks.
\quellseite{359}
\chapter{Das 22. Kapitel}

\section{Vorrede}

Die heilige Biblia muß neben und bei dem Schwerte sein, soll ein Regiment glücklich und wohl von statten gehen; denn die Obrigkeit soll nicht allein Hüter der andern, sondern auch der ersten Tafel sein, und soll nicht allein ihrer lieben getreuen Unterthanen Leib und Gut beschirmen, sondern auch daran sein, daß die Religion und der rechte Gottesdienst wohl bestellt sei. Denn wie eine Laterne ein Licht bewahret und erhält, daß es nicht liederlich kann ausgeblasen oder ausgelöscht werden, und, wenn auch das Licht in der Laterne nicht brennet, so ist sie zu lauter nichts nütze: also sollen auch christliche Regenten das helle Gnadenlicht Christum und sein göttliches Wort lassen lehren und predigen, daß der lebendige persönliche Gnadenthron in der ganzen Welt herum geführet werde, und sollen darob sein, daß der Fürst der Finsterniß, der Teufel, solch Gnadenlicht, weil es ihm in die Augen scheinet und unleidlich ist, nicht auslösche; ohne dies ist ihre Regierung auch nicht allerdings nütze. Dies lehret auch das 22. Kapitel.

\section[Die dritthalb Stämme richteten im Heimkehren einen Altar]{Die dritthalb Stämme richteten im Heimkehren einen Altar auf}

\section{Summarium}

Ist eine schöne nützliche Historia, in welcher die weltliche Obrigkeit ihres fürnehmsten Amts erinnert wird, daß neben dem sie die Unterthanen schützen und ihnen treulich fürstehen soll, sie auch dem Exempel Josua nach, ihr Volk zu Gottes Wort halten, ihnen dasselbige befehlen und fleißig einbilden soll, V. 5. Denn ob wohl das Predigtamt ein unterschiedliches Amt ist von dem weltlichen Regiment, so soll doch die weltliche Obrigkeit das Wort Gottes mit Fleiß helfen fördern, und die Strafe mit Ernst gehen lassen wider alles so dem Wort entgegen und ärgerlich ist. Wo Gottes Name also von der weltlichen Obrigkeit geehrt und dem Aergerniß gewehret wird, da will Gott die Leute wiederum mit Glück und Segen ehren. Wiederum, wo weltliche Obrigkeit Gottes Ehre nicht achtet, und nimmt sich um die Kirche nicht an, da will Gott des Regiments sich auch nicht annehmen, und muß endlich alles zu Trümmern gehen, und kann weder Gewalt noch Reichthum helfen.
% Scan page 360: No non-biblical apparatus, gloss, or structural title/synopsis retained in OCR blocks.
\quellseite{361}
\section{Summarium}

Sie ist eine schöne Historia, wie die zehn Stämme besorgen, es möchten die andern jenseit des Jordan Israel sich zur Abgötterei geben, daß sie ohne Befehl einen Altar in Gilead gebauet hatten; aber da sie höreten, wie solcher Altar nicht zum Opfer gebauet wäre, sondern zum Zeugniß, daß sie auch zu GOttes Volk gehörten, ließen sie es ihnen gefallen, V. 30. Dies Exempel zeigt an, mit was Fleiß und Ernst weltliche Obrigkeit darauf sehen solle, daß allem Aergerniß gewehret werde.

\section{Votum}

Der treue, barmherzige GOtt lasse ihm unsere liebe Obrigkeit, hohen und niedrigen Standes, befohlen sein, und helfe, daß ihre Regierung auch ferner zur Erhaltung seines Worts diene, und daß allem Aergerniß gewehret werde, zu Lob und Preis seinem Namen, und den Unterthanen zur Wohlfahrt und Seligkeit, Amen.

\chapter{Das 23. Kapitel}

\section{Vorrede}

Wer nicht Gottes offenbarem Wort folget in all seinem Thun und Lassen, der gedenke nicht, daß er glücklich sein und Gottes Segen zu Leibes und der Seelen Nothdurft haben und behalten werde; denn nach dem Unglauben ist Gott keiner Sünde feinder, strafet auch keine härter und peinlicher, als wenn man wider sein Wort, darin er seinen ewigen und endlichen Willen offenbaret hat, muthwillig handelt. Dargegen daß in den schweren, betrübten, gefährlichen Zeiten Gott uns auch segne, so lasset uns nach seinem Wort und Geboten leben, und laßt uns täglich herzlich beten mit dem königlichen Propheten David, Psalm 119, 33.: Zeige mir HErr den Weg deiner Rechte, daß ich sie bewahre bis ans Ende. Unterweise mich, daß ich bewahre dein Gesetz und halte es von ganzem Herzen. Führe mich auf dem Steige deiner Gebote, denn ich habe Lust dazu. Dies erinnert auch das 23. Kapitel.

\section{Josua hält einen Landtag, das Volk zu vermahnen}
\quellseite{362}
\section{Summarium}

Hier ist eine sehr schöne Vermahnung, durch welche Josua als ein frommer und gottesfürchtiger Fürst, vor seinem Tod sein Volk vermahnet, daß sie sich nach dem Wort und Willen Gottes halten sollen, so werde Gott Glück und Heil ihnen widerfahren lassen, V. 6. Thun sie es aber nicht, sondern weichen ab von Gott und seinem Wort auf der Heiden abgöttisches Wesen, so werde Gott mit allerlei Unglück sie überschütten, bis sie endlich gar umkommen.

Solches Exempel dieses gottseligen Fürsten sollte alle Obrigkeit mit Fleiß bedenken, und ihm folgen, auf daß die Unterthanen durch ihrer Obrigkeit Exempel und ernstlichen Befehl lernen sich für allem Aergerniß hüten, und gottselig leben.

\section{Votum}

Helffe GOtt, der Herrscher aller Welt, daß Obrigkeit und Unterthanen seinem heiligen Wort folgen, auf daß beide seinen Segen haben und behalten an Land und Leuten, an Gut und Nahrung, an Leib und Seele bis zum ewigen Leben, Amen.

\chapter{Das 24. Capitel}

\section{Vorrede}

Gott pfleget die Undankbarkeit zu strafen, daß er von den undankbaren seine Gaben wiederum hinweg nimmt, wie er denn allen denen gethan hat, welchen er sein Wort gegeben, und sie es verachtet haben, und will solches noch thun. Wie davon lehret St. Paulus 2 Thess. 1, 10., da er also sagt: Dafür, daß sie die Liebe zur Wahrheit nicht haben angenommen, daß sie selig würden; darum wird ihnen Gott kräftige Irrthümer senden, daß sie glauben der Lüge: also, wenn man Gott nicht dankbar ist für den Frieden, so sendet er Krieg; wenn man nicht dankbar ist für die Gesundheit, so sendet er Krankheit, und will es dazu bei diesen zeitlichen Strafen nicht bleiben lassen, sondern er will die Undankbarkeit auch mit dem höllischen Feuer strafen. Derowegen weil wir bekennen müssen, daß bei diesen geschwinden Zeiten dennoch Gott dieser Stadt und Land viel geistliche und leibliche große Wohlthaten erzeiget hat, so laßt uns solche erkennen, und Gott mit Herz, Mund und gottseligem Leben dafür danken, auf daß er unserer Undankbarkeit wegen auch nicht vollends seine Gnade von uns wende, und denn zeitliche und ewige Strafen folgen. Darauf höret ferner an das 24. und letzte Kapitel.

\section{Letzter Landtag Josuas. Sein und Eleasars Tod}
\quellseite{363}
\section{Summarium}

Ist eine letzte Predigt, welche Josua der gottselige und fromme Fürst seinem Volk thut, daß sie bedenken wollen, was für große Wohlthaten ihnen Gott erzeiget habe, und sich hüten, daß sie nicht undankbar sein, und von solchem Gott und seinem Wort nicht abweichen, V. 2. Solcher Befehl ist trefflich nöthig, sonderlich wo neue Regenten angehen, denn da folget gemeiniglich Aenderung. Es geschieht aber solche Aenderung selten zur Besserung; wo alsdenn das zuschlägt, da man Gottesfurcht und Wort aus den Augen läßt, da muß Unglück mit Haufen folgen.
\quellseite{364}
\section{Summarium}

Ist ein fein Exempel, wie der gottseelige und fromme Fürst Josua, nachdem das Volk sich erboten, sie wollten von Gott nicht abweichen, sie auf ein neues mit einem Eide verpflichtete, auf daß sie nach seinem Tode solcher Verpflichtung gedächten, und desto weniger abweichen sollten. Solchem Exempel sollten heutiges Tages und allezeit die weltlichen Regenten folgen; aber leider es geschieht bei Wenigen. Deren aber findet man viel, die ihre Unterthanen wider Gott, wider alle Rechte und Billigkeit beleidigen, daß sie andere Papisten-Abgötterei halten, und das Evangelium nicht allein nicht annehmen, sondern auch verfolgen sollen. Und zu solchem treibet und bringet der Papst und seine Bischöfe; diese werden ihr Urtheil finden. Aber fromme gottseelige Obrigkeit soll dies Exempel wohl merken, und alle Mittel und Wege brauchen, daß Gottes Wort rein und lauter, und der Gottesdienst ungefälscht bleibe.

\section{Votum}

Also hat eure Liebe das Buch Josua gehöret, in welchem wir sehen, wie GOtt für sein Volk gestritten und die Heiden um ihrer Sünden willen gestrafet und vertilget hat; uns zum Exempel, daß wir uns bessern, und nicht auch dergleichen Strafe verdienen. Das verleihe uns unser lieber HErr JESus Christus, Amen.

\part{Das Buch der Richter}

\chapter{Das 1. Kapitel}

\section{Vorrede}

Das Buch der Richter ist ein Exempelbuch der Abwechselung, beides der Strafen und Erlösung der Kirchen Gottes. Denn wie nach Gottes Ordnung Nacht und Tag, Winter und Sommer sich abwechseln und verändern: also ist aus sonderem weisen Rath Gottes, der die Seinen wunderbarlich führet, und nichts so ernstlich und begierig sucht, als ihrer aller ewiges Heil und Seligkeit, eine stetige Abwechselung in der christlichen Kirchen, des Leides und der Freude, des Kreuzes und Trostes. Ins sonderheit aber sollen alle fromme Herzen ihnen wohl einbilden die Geschichte von Gideon und Simson, weil in diesen der HErr Christus und seine Kirche schön abgebildet sind. Denn als Esaias der Prophet im 9. Kapitel des HErrn Christi Wohlthaten beschreibt, so deutet er auf die Historia Gideons, und sagt: Denn du hast das Joch ihrer Last, und Ruthe ihrer Schulter, und den Stecken ihres Treibers zerbrochen, wie zur Zeit Midian. So wird unser HErr Christus ein Nazarener genannt, Math. 2. Welcher Name auch eigentlich zeiget auf die Geschichte von Simson. Darauf höret nun an das 1. Kapitel.

\section{Der Stamm Juda führet an Josuas Statt den Krieg}
% Scan page 365: No non-biblical apparatus, gloss, or structural title/synopsis retained in OCR blocks.
\quellseite{366}
\section{Summarium}

Hierin werden erstlich etliche Historien wiederholet, so zur Zeit Josua sich begeben haben. Darnach wird eine sonderliche Historia angezeigt, wie Gott den Tyrannen AdoniBesek habe gestraft, V. 5. Denn Gott, ob er wohl gebeut, daß weltliche Obrigkeit das Uebel mit Ernst soll strafen, so will er dennoch alle Tyrannei verboten haben. Wie es denn eine sonderliche Tyrannei an dem AdoniBesek gewesen ist, daß er siebenzig Könige an ihrem Leibe unbillig verstümmelt, und sie wie die Hunde das Brod unter dem Tisch hat lassen auflesen. Darum ward ihm vergolten, wie ers an andern verdienet hatte, V. 7. Solch Exempel soll weltliche Obrigkeit bedenken, und gleich wie sie auch nicht zu gelind wider öffentliches Aergerniß soll sein, also soll sie auch in der Strafe nicht tyrannisch sein, und den Zorn einfangen, und Maß halten. Denn es ist bald geschehen, daß man zu weit greift, wo man den Zorn läßt walten.

\section{Votum}

Der treue GOtt lindere uns die wohlverdienten Strafen, und regiere die Obrigkeit, und uns Unterthanen, daß die Obrigkeit alle Laster ernstlich strafe, und daß wir Unterthanen, gegen ihn und alle Menschen, hohen und niedern Standes, christlich leben, Amen.

\chapter{Das 2. Kapitel}

\section{Vorrede}

Auf gleiche Sünden folgen gleiche Strafen. Das lehret Gottes Wort, auch bezeugets die Erfahrung. Dennoch sind die Sünden jetzt am gemeinsten, welche Gott allzeit am greulichsten gestraft hat, als da sind Abgötterei, Hoffart, Wucher, Geiz, Unzucht, Büberei und dergleichen. Derwegen, ach ein jeder bedenke sich, wie er lebe, und ob sein Wandel Gott gefalle, ihm und andern nützlich oder schädlich sei. Der Fromme fahre fort in der Frömmigkeit nach Gottes Wort. Wer in Sünden lebet, der thue Buße, lasse von Sünden, und thue Gutes, und lebe also auch nach Gottes Gebot, daß er ihm und denen, derer er ja schonen soll, nicht Herzeleid und Strafe zurichte. Dies soll auch erinnern das 2. Kapitel.

\section{Der Kinder Israel Trägheit, Abfall, Strafe und Erledigung}

a Engel) der Priester Pinehas.

b Bochim heißt die Weinenden.
\quellseite{367}
\section{Summarium}

Ist gleich wie eine Vorrede über dies Buch, in welchem fürnehmlich wird angezeigt, wie die Kinder Israel auf der Heiden Abgötterei gerathen, und von Gott und seinem Wort sind abgewichen, V. 2. Und wie Gott solche Sünde habe damit gestrafet, daß sie ihren Feinden nicht mehr widerstehen möchten, sondern die Hand Gottes wider sie zum Unglück gewesen sei, wo sie hinaus wollten; auf daß man lerne in Gottes Furcht und Willen wandeln, sonst soll es uns auch also gehen.

Sonderlich aber ist zu merken, daß der Text spricht: Das Volk dient dem HErrn, so lang Josua lebte, und die Aeltesten nach ihm, V. 7. Denn wo fromme Regenten abgehen, die das Volk zum Wort und Gottesdienst halten, ist es leichtlich geschehen, daß das Wort sich wieder verliere. Wie wir denn heutiges Tages bei dieser jungen armen Welt auch besorgen müssen, wenn Gott diese wird hinweg nehmen, durch welche er uns sein Wort hat offenbaret.

\section{Votum}

Der barmherzige GOtt lindere die wohlverdienten Strafen, behüte uns für Abgötterei, und erhalte dieser Stadt und dem ganzen Lande sein heilig Wort bis an jüngsten Tag, und bewahre uns in wahrem Glauben und gottseligem Wandel, bis zum seligen Tod und Ende, Amen.

\chapter{Das 3. Capitel}

\section{Vorrede}

Es ist wahr, Gott erhöret die Sünder nicht, Joh. 9, 31. Aber dies soll verstanden werden von unbußfertigen Sündern, welche in offenbaren Sünden leben, und weiter zu sündigen gedenken. Denn also sagt Salomon Sprüchw. 28, 9.: Wer seine Ohren abwendet zu hören das Gesetz, deß Gebet ist ein Greuel. Aber die Sünder, die Buße thun, und in wahrem Glauben beten, die werden erhöret. Davon sagt Gott also 2 Chron. 7, 14.: Wenn sie beten, und mein Angesicht suchen, und sich von ihren bösen Wegen bekehren werden, so will ich vom
\quellseite{368}
Himmel hören, und ihnen ihre Sünde vergeben, und ihr Land heilen. So helfe der allmächtige, gütige Gott, daß wir allezeit aus bußfertigen Herzen thun unser Gebet, und mit David sprechen Psalm 5, 8.: Ich aber, HErr, will in dein Haus gehen, auf deine große Güte, und anbeten gegen deinem heiligen Tempel in deiner Furcht, das ist, mit einem bußfertigen, gläubigen Herzen. Darauf höret ferner an das 3. Kapitel.

\section{Athniel, Ehud, Samgar, und ihre Thaten}

\section{Summarium}

Dies ist eine Historia, wie Gott etliche Heiden habe lassen überbleiben, auf daß er sein Volk versuchte, ob sie auch an seinen Geboten halten, und nicht der Heiden Exempel nach Abgötterei treiben wollten, V. 1. Aber sie haben sich das Aergerniß lassen übergehen, und gesündigt, V. 6. Darum hat Gott der Syrer König wohl acht Jahr über sie herrschen lassen, V. 8. Solches soll uns reizen, daß wir Gott fürchten, und in seinem Willen und nach seinem Wort wandeln. Sintemal er die Sünde durch Krieg und fremde Herrschaft pflegt zu strafen. Wiederum ist es tröstlich, daß Gott dennoch solcher Sünder Gebet erhöret, wenn sie um Hülfe schreien, und läßt ihnen Hülfe widerfahren.
\quellseite{369}
a Die Rathstube, ubi ordinatae erant sedes.

\section{Summarium}

Ist wiederum ein Exempel, wie Gott sein Volk um der Sünde willen mit Krieg und fremder Herrschaft plaget, aber hernach, da sie um Hülfe schreien, wieder durch einen frommen Fürsten errettet habe, V. 15. Daß nun Ehud den König der Moabiter heimlich mit Eisen ersticht, V. 21., ist nicht dahin zu deuten, als sollten solche Exempel Ausführer und Mörder billigen, sondern man muß auf den Befehl, das Amt und den Beruf Ehud sehen, der war von Gott dazu berufen, daß Israel durch ihn von den Moabitern sollte errettet werden, und den König erwürgen; wie er denn zum Könige sagt, ich habe Gottes Wort an dich, V. 20., das ist, ich habe einen Befehl von Gott, der dich belanget. Wo nun solcher sonderlicher, gewisser und ausgedrückter Befehl Gottes nicht ist, da soll man sich aller Gewalt enthalten, und auch das Unrecht geduldig leiden, V. 25., bis Gott selbst durch andere Mittel von den Tyrannen erledigt und hilft.

\section{Votum}

Dem allmächtigen und barmherzigen GOtt sei Lob und Dank, der die wohlverdienten Strafen lindert; der wolle ferner unser Gebet erhören, uns in allen Nöthen beistehen, und unter dem Schutz unserer lieben Obrigkeit uns segnen und behüten, Amen.

\chapter{Das 4. Kapitel}

\section{Vorrede}

Im Psalm 50, 25. spricht der HErr: Rufe mich an in der Noth, so will ich dich erretten, so sollst du mich preisen. Hiermit lehret Gott, da jemand gesündiget hat, daß er sich seine Sünde nicht soll von Gott abschrecken lassen, sondern soll zu ihm bringen ein zerknirschtes Herz, und soll sich des HErrn Christ in wahrem Glauben trösten, und einen guten Vorsatz haben, sein Leben zu bessern, und soll also ihn anrufen, so wollte er erhören und helfen. Dies lehret auch St. Paulus 2 Timoth. 2, 19. und spricht: Es trete ab von Ungerechtigkeit, der den Namen Christi anrufet. Und Salomon in seinen Sprüchen 15, 8. sagt also: Des Gottlosen Opfer ist dem HErrn ein Greuel, aber das Gebet der Frommen ist ihm angenehm. Deswegen so laßt uns Buße thun, und Gott im Geist und in der Wahrheit anrufen. Es ist gewiß, wir werden entweder Abwendung oder ja Linderung der gegenwärtigen Strafen erlangen. Darauf hört auch an das 4. Kapitel.

\section{Debora, Barak und Jael's Heldenmuth}
\quellseite{370}
\section{Summarium}

Ist eine schöne Historia, wie Gott das Volk um ihrer Sünde willen strafet, und hernach, da sie um Hülfe schreien, erhöret, und ihnen hilft durch die Prophetin Debora, dazumal Richterin im Volk Israel, V. 4.; auf daß man sehe, wie Gottes Hülfe nicht an großer Macht und Gewalt liegt. Er kann ebenso wohl durch ein Weib erretten und helfen, als durch einen Mann, wo man sich nur seines Willens hält, und nicht Ursach zur Strafe gibt. Denn es liegt alles an dem, ob Gott den Feinden das Herz nehme oder nicht.

\section{Summarium}

Ist ein fein Exempel, wie Gott den Tyrannen eine Weile pfleget zuzusehen; wenn aber ihr Stündlein kommt, müssen sie einer leichten Ursach halben zu Boden gehen. Wie denn der tapfere Kriegsmann Sisera von einem Weib erwürget wird, unangesehen, daß zwischen ihrem Hause und dem König Jabin Friede war.
\quellseite{371}
Denn an dem war ihr, als einer Regentin, die Gottes Befehl dazu hatte, mehr gelegen, daß ihr Volk errettet würde, denn an der Privat-Verwandtnis, so zwischen ihnen war.

\section{Votum}

Der treue, gütige GOtt stehe auch heute seiner Kirche bei, steure und wehre allen deren Verfolgern und Tyrannen, und verleihe unserer lieben Obrigkeit ferner ruhige und glückliche Regierung, den armen Unterthanen zu seliger Wohlfahrt, Amen.

\chapter{Das 5. Capitel}

\section{Vorrede}

Wie wir nachlässig zum Gebet sind, also sind wir viel nachlässiger zur Danksagung. Die Undankbarkeit aber kann Gott nicht leiden, wenn dieselbige schon ein Mensch dem andern beweiset, wie wir lesen 1 Mos. 31, 12., da er die Undankbarkeit Laban gegen Jakob strafet; und zuvor 1 Mos. 16, 9., da er auch die Undankbarkeit Hagar gegen Sara strafet. Denn es stehet geschrieben in Sprüchen Salomonis 17, 13.: Wer Gutes mit Bösem vergilt, von deß Hause wird Böses nicht lassen. Viel weniger wird nun Gott leiden, daß man ihm für seine Wohlthaten undankbar ist. Derwegen, Jünglinge und Jungfrauen, Alte mit den Jungen, sollen allzeit loben den Namen des HErrn, Psalm 148, 12. Und sollen ohn Unterlaß beten, und in allen Dingen danksagen, 1 Thessal. 5, 17. Darauf höret an das 5. Capitel.

\section{Deboras und Baraks Triumphlied}

a Dies Lied will so viel sagen, daß Gott habe den Sisera geschlagen durch die geringsten Leute in Israel, daß die Geringen auch einmal hoch und groß worden sind, da die großen Geschlechter Israel stille saßen, und sie verblieben in Nöthen; daß ist das Neue, das der HErr ermahlet hat, da sind die Bauern Israel prächtig, und auch Herren worden.

b Krumme) das ist, es war kein Regiment noch Ordnung im Lande.

c Lobet) das ist, ihr Herren, Richter und gemeine Leute.

d Schöpfern) das ist, da die Schützen Sisera schrieen für Noth am Wasser Kison, da man pfleget zu schöpfen, da half Gott seinem Bauern, und ließ das Recht gehen.

e Wurzel) das ist, Josua war der erste Fürst aus dem Stamm Ephraim, der schlug Amalek, und nach ihm die andern, bis daß Sebulon auch einmal einen Josua überkommen hat wider Sisera. Und merke, sie nennet die Fürsten Regierer, die mit der Schreibfedern streiten, das ist, sie gewinnen (geistlich davon zu reden) mehr durch den Glauben an Gottes Wort, denn mit dem Schwert.

f Hürden) das ist, du bleibest daheim, ob du wohl hörtest das arme Häuflein zu Felde blasen, und hattest doch nahe zu ihm.
\quellseite{372}
\section{Summarium}

Die ist ein Danklied, daß Gott seinem Volk wider die Cananiter geholfen hat. Denn bei den Frommen sollen diese zwei Stücke auf einander gehen: daß man erstlich erkenne, wie Gott alles aus lauter Güte und Barmherzigkeit uns widerfahren lasse; das andere, daß man auch dafür dankbar sei, und solcher Wohlthat nicht vergesse.

Am Ende ist ein feiner Spruch, daß die Prophetin wünschet, daß alle Feinde Gottes sollen umkommen, wie Sisera, aber die Gott lieben, sollen aufgehen, wie die Sonne in ihrer Macht, welche niemand dämpfen kann, sondern muß sie frei gehen und leuchten lassen, V. 13. Solcher Wunsch widerfähret durchaus der christlichen Kirchen, die hat die Anfechtung und Hinderniß in der Welt sehr gewaltig, daß es scheinet, als sollte es alles zu Boden gehen. Aber wie Debora hie bittet, dringet sie doch immer fort wie die liebe Sonne, so ihr Licht nicht einfangen noch aufhalten läßt, wie Christus sagt, daß auch der Höllen Pforten sie nicht bewältigen sollen.

\section{Votum}

Lob, Ehre, Preis und Dank sei GOtt für alle seine Wohlthaten; der wolle ferner uns und die ganze Christenheit unter der vorgefallenen wunderlichen Zeit und Veränderungen segnen, behüten und erhalten, daß bis an den jüngsten Tag bleiben, die ihn erkennen und preisen, Amen.

\chapter{Das 6. Kapitel}

\section{Vorrede}

Wer Glück in seinem ehrlichen Beruf, Amt und Stande haben will, der muß dreierlei in Acht nehmen. Erstlich soll er gewiß sein, daß sein Beruf Gott gefalle; dann soll er sich in seinem Beruf alles möglichen Gehorsams gegen Gottes Wort befleißen; und zum dritten soll er Gott um glücklichen Fortgang seines Berufs täglich und herzlich anrufen. Deß ist ein Exempel der theure Fürst Gideon, von dem das 6. Kapitel.

\section{Gideon zum Richter berufen}
\quellseite{373}
\section{Summarium}

Hie fähet an die Historia des theuern und gottseeligen Fürsten Gideonis, welchen Gott dazu berufet, daß er ein Fürst sein und dem Volk Israel wider die Midianiter helfen soll. In dieser Historia sehen wir erstlich, wie Gott die Sünden durch Krieg und Unfrieden pfleget zu strafen; aber dennoch die Strafe abwendet, wenn man Hülfe bei ihm suchet. Daß also durchaus solche Historien eitel Exempel der Buße und des Glaubens sind, daß wir uns für Sünden sollen hüten, weil sie Gott mit strafen, und wiederum auf Gottes Güte sehen, und uns derselben trösten, wenn wir die Strafe empfunden und tragen; sintemal er hie so gnädig erhöret, da man ihn bittet, und sich Hülfe zu ihm versiehet. Darum ist sonderlich zu merken an Gideon, da der Engel ihn grüßet: Der HErr mit dir, daß ers beredet und spricht: Ist der HErr mit uns, warum gehet es denn so übel? V. 13.
\quellseite{374}
Denn also pflegt die Vernunft zu urtheilen, wenn es übel gehet, daß Gott gar sey von uns gewichen; so doch Gott das Unglück uns darum auferlegt, daß wir uns zu ihm halten, und bessern, und seine Hülfe suchen sollen.

\section{Summarium}

Ist ein sehr schön Exempel: ehe Gott seinem Volk durch Gideon hilft, befehlet er, daß er zuvor den Altar Baal abbrechen, und den Hain oder Forst umhauen soll. Denn wo Glück und Heil soll sein, da muß am ersten durch die Obrigkeit Abgötterei und andere Sünde abgeschafft sein.

Daß nun Gideon zwei Zeichen mit dem Fell fordert, V. 36., geschieht alles darum, daß er seines Berufs desto gewisser sei, und mit größerem Muth die Sachen angreifen möge; wie denn die Historia hernach anzeigen wird, wie ihn Gott zu einem sonderbaren hohen Werck hat brauchen wollen, daß es sehr demüthig war, daß er zuvor seines Berufs gewiß wäre.

\section{Votum}

GOtt sei uns gnädig, erhöre unser Gebet und helfe, daß aller christlichen Obrigkeit Regierung diene zur Heiligung seines Namens, daß sein Reich zu uns komme, und daß sein Wille geschehe, wie im Himmel, also auch auf Erden, von uns und allen Menschen, Amen.

\chapter{Das 7. Capitel}

\section{Vorrede}

Des Glaubens Natur und Art ist, daß er sich fröhlich erwäget, die Dinge zu glauben, die er nicht fühlet, noch mit allen seinen Kräften inwendig und auswendig ermessen kann. Welche
\quellseite{375}
Menschen nun einen solchen Glauben haben, und in dem an Christum Jesum festhalten, die macht er zu Herren, wie er selbst ist, über Sünde, Tod, Teufel, Welt und alle Dinge. Derwegen trachten wir billig nach dem wahren Glauben und seufzen: O Herr, Gott, heiliger Geist, stärk unsern Glauben immerdar, daß wir an Christo bleiben, dem treuen Heiland, und ihm aus ganzer Macht allzeit vertrauen, Amen. Dies erinnert auch das 7. Kapitel.

\section{Gideon erhält den Sieg wider die Midianiter}

\section{Summarium}

Ist eine herrliche Historia, wie die Midianiter durch Gideon erlegt sind. Denn hie stehet man,
\quellseite{376}
wie einen großen Glauben der Fürst auf Gott und seine Zusagung gehabt hat, daß er nach dem Befehl Gottes das Volk wieder heimziehen heißet, und nicht mehr bei sich behält denn dreihundert Mann, so doch der Feinde eine unzählige Menge war, V. 3.; auf daß wir lernen, wie alles an Gott und seiner Hülfe, und nicht mehr an Menschen Vernunft, Weisheit, noch Macht liege; aber solches glaubt die Vernunft nicht.

\section{Summarium}

Ist die Historia, wie die Midianiter wunderbarlich ohne allen Schwerthschlag erlegt sind. Solche Historia zeucht der Prophet Es. 9, 7. als eine Figur, auf die Erlösung, so uns durch unsern HErrn Christum widerfahren ist, da wir von dem geistlichen Midian, das ist von der Sünde, Tod und Gesetze, und des Teufels Gewalt erlöset sind, ohne unser Zuthun, allein durch Christus Sterben und Auferstehung.

Mit solcher Auslegung reimet sich auch die heimliche Deutung. Die Midianiter sind, wie gesagt, der Christen ärgste Feinde, nämlich, die Sünde, der Tod, Welt, Teufel und Fleisch. Solche Midianiter lassen sich mit dem Schwert nicht schlagen, sondern das muß die Rüstung sein, daß wir in der rechten Hand haben die Posaunen, das ist das Wort Gottes, und in der linken die Lampen des Glaubens, welche wir tragen im irdischen Gefäß, wie St. Paulus sagt. Desselbige Gefäß muß zerschlagen werden, das ist, Fleisch und Blut muß gedämpft und täglich getödtet werden, bis wir endlich gar sterben; so gehet es denn, daß unsere Feinde überwunden und geschlagen, wir aber durch Christum selig werden, Amen.

\section{Votum}

Wir danken dem HErrn Jesu Christo für seine Victori und Sieg wider seine Feinde. Der wolle unsern schwachen Glauben in uns stärken, und helfen, daß wir auch überwinden alle unsere geistlichen und leiblichen Feinde, im Leben und Sterben, Amen.

\chapter{Das 8. Kapitel}

\section{Vorrede}

Der heilige Apostel St. Jacob sagt in seiner Epistel am 1. Kapitel V. 11.: Die Sonne gehet auf mit der Hitze, und das Gras verwelket, und die Blume fällt ab, und sein schöne Gestalt verdirbet. Also hat ein jeder Mensch sein Kreuz, das von Tage zu Tage des Leibes und aller Gliedmaßen Kräfte schwächet, bis endlich folget der bittere Tod, der keinen Unterschied hält unter den Menschen, und derhalben sobald ein Junges würget, als ein Altes, dem niemand Widerstand thun kann, der die Menschen oft plötzlich überfällt und hinreißt, der auch niemand zuvor lässet wissen, wenn, wie oder wo er sterben werde. Demgegen wollen wir Buße thun, demüthig sein, Gottes Wort hören, wachen und beten, und das neben christlich leben, auf daß wir darauf auch fröhlich und selig sterben in Christo Jesu, unserm HErrn, Amen. Darauf höret weiter an das 8. Kapitel.

\section[Die übergebliebenen Midianiter werden geschlagen, und Sebah und]{Die übergebliebenen Midianiter werden geschlagen, und Sebah und Zalmunna gefangen}
\quellseite{377}
a Das ist, erfahren, inne werden.

\section{Summarium}

Ist ein fein Exempel, wie der fromme Fürst Gideon sich nicht erhebt um des großen Sieges willen, welchen ihm Gott gegeben hatte, sondern demüthiget sich, und durch solche Demuth erhält er Friede, da sonst durch Hoffart und Stolz, wie allwege geschicht, Unfriede wäre erreget worden.
\quellseite{378}
\section{Summarium}

Ist eine sonderliche Historia, wie oft etwas, daß man guter Meinung, auch aus guten Ursachen thut, übel geräth und groß Unglück verursachet, wie hie Gideon fordert die güldenen Ringe von den Kriegsleuten, und machet einen Leibrock daraus, V. 27., auf daß es solt ein Denkzeichen sein des herrlichen Sieges, welchen ihm Gott gegeben hat. Aber es gerieth solcher Anschlag übel. Denn sie mißbrauchten hernach solches Leibrocks zu Abgötterei, und hielten ihn dem gleich, welchen Gott befohlen hatte zu machen und der Hohepriester anlegen mußte. Es ist aber solcher Abgötterei halben, wie folgen wird, das ganze Volk hart gestraft worden.

Darum soll jedermann darauf sehen, daß er nichts fürnehme, noch anfange, er habe denn Gottes Wort und Befehl; wer dabei bleibt, der darf sich keines Mißbrauchs oder einiger Sünde besorgen.

\section{Votum}

Helfe Gott, daß sein Wort eine Leuchte sei unserer Füße, und ein Wegzeiger in unserm Beruf, und regiere uns, daß wir uns für allerley Abgötterei und andern Sünden und Lastern hüten, bis an unser selig Ende, Amen.

\chapter{Das 9. Capitel}

\section{Vorrede}

Niemand soll seines Nutzes wegen wider Gott und seinen Nächsten handeln, denn was dawider geschieht, obs schon eine Zeit lang einen glücklichen Fortgang hat, so wirds doch endlich gestraft. Derwegen nehme ein jeglicher wahr in seinem Berufe seines Nächsten, und schade nicht, sondern diene ihm, so sammlet er ihm und den Seinem Gottes Segen. Dieses lehrt St. Paulus 1 Corinth. 13., da er sagt: Die Liebe suchet nicht das Ihre, das ist: sie trachtet nicht auf ihr eigen Gut, Nutz, Ehr und Gemach, sondern trachtet darauf, wie sie des Nächsten Nutz, Ehre, Leib und Leben fördere. Ach, das helfe uns der treue Gott. Denn weil wir von Natur alle eigennützig sind, so ist wohl zu achten, daß wir die christliche, brüderliche Liebe nicht von uns selbst haben, sondern Gott muß sie uns geben durch den heiligen Geist. Ferner höret an den ersten Theil des 9. Capitels.

\section{Der Brudermörder Abimelech kommt um}
\quellseite{379}
a Es bestehet nicht, es ist ein Schuchal.

\section{Summarium}

Ist eine Historia, wie der Tyrann Abimelech sich des Regiments unterstehet und würget siebenzig Söhne des Fürsten Gideonis, auf daß das Regiment ihm allein bliebe, V. 5. Ob nun wohl solche Strafe Gideon verdienet habe, damit, daß er mit dem Leibrock, den er machen ließ, zur Abgötterei Ursache gab, lasse ich bleiben. Gleichwohl sagt der Text Jud. 8, 27. mit klaren Worten, daß solche That, die doch aus guter Meinung fürgenommen war, nicht allein Gideon und seinem Hause zum Aergerniß sei geraten, sondern daß auch das ganze Israel damit Abgötterei getrieben habe. Dadurch wir ja treulich gewarnet werden, daß wir bei Leibe nichts fürnehmen oder anfangen sollen aus eigener Wahl, guter Meinung und Andacht, Gott damit zu dienen und ehren, es habe einen Schein der Heiligkeit so groß und gut es kann. Geschichts ohne Gottes Wort und Befehl, so gibt es Ursache zu allerlei Irrthum und Abgötterei, wenns noch einst ein heiliger Patriarch Gideon aus sonderlicher Andacht angerichtet hätte. O wie sicher, frech und durstig hat hinein gefahren der höllische Drache zu Rom, Gott mache solches Lästernds schier ein Ende, Amen.

\section{Der andere Theil des 9. Kapitels}

\section{Vorrede}

Um die Sünde ists so ein arg Ding, daß dieselbigen von der Erde gen Himmel schreien, und für Gottes Angesicht gleich treten, und die Rache und Strafe erfordern, daß er endlich seinen Zorn darauf muß lassen ergehen, wo man sich nicht durch wahre Buße zu ihm bekehret. Die schreienden Sünden aber werden von Alten diese genennet: Todtschlag, Wucher, Unterdrückung der Wittwen und Waisen, Sodomie, und verdienten Lohn fürhalten. Derwegen wollen wir ablassen von unsern Sünden, und die wir leider gethan, wollen wir Gott abbitten im Namen und auf die Vorbitte unsers Herrn Jesu Christi, und wollen Gott und unserm Nächsten dienen im Glauben, mit wahrer Liebe, so wird es Gott eine Freude sein, und uns ein Nutzen, zeitlich und ewig. Dazu höret an den andern Theil des 9. Kapitels.
% Scan page 380: No non-biblical apparatus, gloss, or structural title/synopsis retained in OCR blocks.
\quellseite{381}
\section{Summarium}

Ist eine Historia, wie Gott eine Weile den Tyrannen und bösen Buben zu Sichem, so ihm zu seiner Tyrannei geholfen hatten, zusiehet, aber endlich sie bezahlet, wie sie verdienet haben. Und ist sonderlich hie zu merken, daß der Text mit klaren Worten sagt, daß Gott die tyrannischen Regenten mit Aufruhr strafet; und die, so zum Aufruhr helfen, wie hie Gaal thut, nur Schande und Schaden davon bringen. Solche Exempel dienen alle dazu, daß man lerne Gott fürchten und fromm sein. Sintemal Gott die Sünde nicht will ungestraft lassen, wie man hie siehet an den Sichemitern und ihrer Stadt.

\section{Summarium}

Ist eine Historia, wie, nachdem der Tyrann Abimelech die bösen Buben zu Sichem wohl gestraft hatte, er auch muß herhalten und bezahlen, was er mit dem Mord an seinen Brüdern und sonst verdienet hatte. Und ist sonderlich zu merken, daß gemeiniglich die Tyrannen schändlich und liederlich umkommen. Wie den Abimelech ein Weib zu Tode wirft und sein Knecht ihn vollends erstechen muß, V. 53. Auf daß ja die Sünde endlich nicht ungestraft bleibe, und jedermann empfähe, nach dem er verdienet hat.

\section{Votum}

Regiere uns alle der treue GOtt, daß wir der Sünden ja gram werden, und dagegen allen Tugenden hold und günstig, und daß wir leben in Heiligkeit und Gerechtigkeit, Ihm zu Lob und dem Nächsten zum Frommen, in Kraft des heil. Geistes, Amen.

\chapter{Das 10. Kapitel}

\section{Vorrede}

Gott strafet die Sünde; wie solches unsere Ohren hören und unsre Augen sehen. Viele habens auch gefühlet an ihrem Leibe, viele befindens auch an ihrer Nahrung. Ach der Blindheit und wehe der Sicherheit, daß dennoch viele von wissentlichen, muthwilligen und vorsätzlichen Sünden nicht ablassen. Derwegen wollen wir wachen und beten, Gott wolle uns in wahrer Buße erhalten, auf daß wir mit unsern Sünden nicht uns noch andern schädlich sein. Denn der Knecht, der seines Herrn Willen weiß, und hat sich nicht bereitet, auch nicht nach seinem Willen gethan, der wird viel Streiche leiden müssen, sagt unser Herr Jesus Christus Luc. 12, 47. Darauf höret nun an das 10. Kapitel.

\section[Israel wird wegen der Abgötterei von den Philistern]{Israel wird wegen der Abgötterei von den Philistern gedrängt}
\quellseite{382}
\section{Summarium}

Es ist eine Historia, wie Gott sein Volk durch die Ammoniter und Philister um ihrer Abgötterei und Sünden willen gestraft habe. Der Text zeiget hie sonderlich an, daß Gott solche Strafe ihnen widerfahren lasse, nicht daß er sie gar verderben, sondern dadurch dahin bringen will, daß sie ihre Sünde erkennen und davon ablassen sollen, V. 14. Daß aber Gott durch den Propheten sagen läßt: Er wolle nicht mehr helfen, sie sollen bei den Götzen Hülfe suchen, welchen sie gedienet haben, V. 15., geschieht darum, auf daß sie ihre Sünde mit größerem Ernst beweinten, und sich desto eher bekehrten. Also gehen gemeiniglich alle Historien dieses Buches dahin, daß man entweder Buße thue und sich bessere, oder gewarte, daß Gott zürnen und heftig durch Krieg und ander Unglück unsere Sünde strafen werde.

\section{Votum}

Der barmherzige GOtt lindere uns die Strafen, so man mit Abgötterei, mit falschem Gottesdienst und dergleichen Sünden verdienet, und helfe, daß wir in seiner Furcht leben, und ihm allzeit dienen nach seinem Wort und Willen, im geistlichen und weltlichen Leben. Amen.

\chapter{Das 11. Kapitel}

\section{Vorrede}

Wohl unserer Stadt und Land, wohl eines jeden Leib und Seele, wohl eines jeden Haus und Nahrung, da wir allezeit und sonderlich in diesen gefährlichen Zeiten fleißig, herzlich und in wahrem Glauben beten. Denn keine Noth ist so groß, und keine Gefahr so grausam, und kein Kreuz so schwer, ein gläubig Gebet kann es entweder abwenden oder lindern. Und der alte Kirchenlehrer Origenes sagt: Eines frommen Christen Gebet ist stärker, und richtet mehr aus, als die gewappnete Hand vieles Kriegsvolks. Denn dieweil das Gebet den Himmel durchbringet, sollte es dann nicht auch auf Erden mächtig sein und alle Feindschaft überwinden? Dies ist fürnehmlich ein Exempel des 11. Kapitels.

\section{Jephthahs Sieg und Gelübde}
% Scan page 383: No non-biblical apparatus, gloss, or structural title/synopsis retained in OCR blocks.
\quellseite{384}
\section{Summarium}

Ist eine Historia, wie Gott sein Volk durch den Fürsten Jephthah errettet habe, da sie sich gebessert hatten, und zu Gott um Hülfe schrieen. Es ist aber ein fein Exempel hie, daß Jephthah der Untreue vergisset, die ihm von den Seinen widerfahren war, und sie derselbigen nicht will entgelten lassen, da sie seiner Hülfe bedurften, V. 7. Denn Gott will, wo Noth vorhanden ist und wir helfen können, daß man da weder Dank noch Undank ansehen soll, sondern thun, was er befohlen hat.

Auch siehet man hie, wie Gottes Urtheil so wunderlich gehet. Jephthah war seinen Stiefbrüdern zum Mittlern nicht gut genug; aber Gott war er gut genug zum Regenten und Heiland, durch den er sein Volk Israel errettet.

a Beuget, oder demüthigest mich; Gott hat mich hoch erhebt durch diesen Sieg, daß ich mein Haupt hoch und fröhlich aufrichte. Aber du beugest mich, daß ich den Kopf muß niederschlagen mit großem Herzeleid, und solche hohe Freude zum tiefen Herzeleid kehren.

b Wie er gelobet hatte) man wil, er habe sie nicht geopfert; aber der Text stehet da klar. So siehet man auch beyde an den Richtern und Königen, daß sie nach großen Thaten haben auch große Thorheit müssen begehen, zu verhüten den leidigen Hochmuth.

\section{Summarium}

Ist eine sonderliche Historia von dem Gelübde Jephthah, welches er unbedacht that, und mit Sünden vollstreckte. Denn man ist kein Gelübde schuldig, das wider Gott und sein Wort ist. Weil nun Jephthah sein Gelübde sonderlich hat thun wollen, Gott zu Dank und Ehre, sollte er je seine Tochter nicht haben würgen lassen. Sintemal Gott den Eltern befohlen hat, daß sie ihre Kinder in Gottesfurcht aufziehen sollen.

Aber in solchen Exempeln sehen wir, daß auch hohe heilige Leute, die große wunderliche Thaten ausgerichtet, auch große Thorheit begangen haben. Es ist keiner je so heilig noch vollkommen gewesen (das Haupt allein aller Heiligen ausgenommen), der nicht oft sich versündiget, und derhalben hat beten müssen (wie der 32. Psalm V. 5. zeiget): Vergib uns unsre Schuld, und Vergebung der Sünden glauben. Wir arme, elende Leute können dem Erzvater Jephthah, der wie man aus der Historia siehet, ein hochgeleuchter Mann gewesen ist, voll Geistes und Glaubens, und doch gleichwohl schmerzlich strauchelt, das Wasser nicht reichen, noch sind wir so sicher und der Sache, wie wir uns dünken lassen, so gewiß, als könnten wir nimmermehr irren, noch fehlen. Es heißt aber: sei nicht stolz, sondern fürchte dich. Und wer sich dünken läßt, er stehe, der mag wohl zusehen, daß er nicht falle.

\section{Votum}

Helffe der treue GOtt, daß wir uns nach seinem Wort gehorsamlich verhalten, fleißig
\quellseite{385}
beten, in unserm Beruf treu sein; auch unserm Nächsten dienen, auf daß er sich auch nach seinem Wort väterlich gegen uns erzeige und verhalte, Amen.

\chapter{Das 12. Kapitel}

\section{Vorrede}

Salomon in seinen Sprüchen sagt im 16. Kapitel V. 18.: Wer zu Grunde gehen soll, der wird zuvor stolz. Hoffart und stolzer Muth kommt vor dem Fall. Und St. Paulus an die Römer im 13. V. 2. spricht: Wer sich wider die Obrigkeit setzet, der widerstrebet Gottes Ordnung: die aber widerstreben, werden über sich ein Urtheil empfangen. Derwegen menge sich niemand unter die Aufrührischen. Denn es wird kein Aufruhr ungestraft bleiben, und auch derer keiner, so sonst der weltlichen Obrigkeit guter Ordnung mutwillig und wissentlich widerstreben. Sondern jedermann sei unterthan der Obrigkeit, die Gewalt über ihn hat. Kehret und warnet abermal Salomon und St. Paulus, Spr. Salomon 24, 21. Röm. 13, 1. Darauf höret auch davon an das 12. Kapitel.

\section{Niederlage der Ephraimiter}

a Id est, qui ad Ephraim fugerunt in periculis pro salute. Non vos eripitis nos, sed nos vos fugitivos, quia sumus regia tribus. Das ist, die zu Ephraim geflohen sind in der Gefahr ihres Lebens. Nicht ihr errettet uns, sondern wir euch Flüchtlinge; denn wir sind von der königlichen Stamm.

b Schibboleth) heißt eine Kornähre, oder auch wohl eine Mundstrafe.

\section{Summarium}

Ist eine Historia, wie Gott die Ephraiter, so unbilliger Weise aus Hoffart wider den Fürsten Jephthah einen Aufruhr erregeten, hart gestraft hat, daß ihrer zwei und vierzig tausend erschlagen sind, V. 6. Also gehet es: wo man Gott nicht fürchtet, strafet er bisweilen Land und Leute mit Krieg und Aufruhr, aber demnach endlich gehet es über den Aufruhrherrn aus, daß sie müssen zu Grund gehen. Denn Gott will weder Hoffart noch Ungehorsam ungestraft lassen.

\section{Votum}

Der treue GOtt helfe, daß wir uns vor diesem und allen andern Sünden und Lastern hüten, und in Demuth und Gehorsam unserer christlichen Obrigkeit allzeit unterthan sein, derselben zur Freude, uns zu Lob und Ehren und andern zu guten Exempeln, Amen.
\quellseite{386}
\chapter{Das 13. Kapitel}

\section{Vorrede}

Viele Eltern sind in den Gedanken, es stehe ihnen frei, wie oder wozu sie ihre Kinder aufziehen. Viele lassen darüber die Hauszucht also fallen, und die Kinder untereinander aufwachsen, wie das arme dumme Vieh. Damit aber laden die Eltern gar schwerlich auf sich Gottes Zorn, Beschwerung des Gewissens, und endlich schwere Strafe und Rechnung. Dies zu verhüten, lehret die heilige Schrift in Sprüchen und Exempeln, was in dem Theil Gott und ihr Amt von ihnen fordert. Sprüche sind, an die Ephes. 6. Kapitel V. 4. sagt St. Paulus: Ihr Väter, reitzet eure Kinder nicht zu Zorn, sondern ziehet sie auf in der Zucht und Vermahnung zum Herrn. Und Salomon sagt in Sprüchen 22, 15. und 29, 17.: Züchtige deinen Sohn, denn Thorheit steckt dem Knaben im Herzen, aber die Ruthe der Zucht wird sie ferne von ihm treiben. Denn Ruthe und Strafe gibt Weisheit. Ein ausbündig Exempel aber hört an aus dem 13. Kapitel.

\section{Simsons Empfängniß und Geburt}

a Es ging wunderlich zu, wie folget, daß der Engel verschwand, und in der Flamme auffuhr.
\quellseite{387}
\section{Summarium}

Sie gehet an die Historia von dem theuren Fürsten und Helden Simson; in welcher man erstlich siehet, wie es eine große Gabe Gottes ist, wo Gott Eheleuten Kinder gibt, V. 3. Darnach ist sonderlich zu merken, wie die Eltern Simsons des Engels wieder begehren, daß sie erfahren, wie sie ihren Sohn nach dem Willen Gottes sollen auferziehen, V. 8. Denn dies will Amts halben den Eltern sonderlich gebühren, daß sie darauf Achtung haben, wie ihre Kinder in Gottes Furcht und Willen recht auferzogen werden. Und dies ist das rechte Erbe, und der höchste Schatz, welchen die Eltern ihren Kindern können lassen. Aber die Welt, wie man siehet, fraget nach solchem Amt am wenigsten. Jedermann denket, wie er seinen Kindern viel Geld und Gut lasse. Aber weil es an rechter Zucht und Furcht Gottes fehlet, sehen wir, daß weder Glück noch Heil bei den Kindern ist, und es oft viel besser wäre, es ließen die Eltern ihren Kindern weder Heller noch Pfennig, so bliebe mancher Unrath und Sünde dahinten. Darum sollen die Eltern dies Exempel des frommen Vaters Manoah fleißig merken.

\section{Votum}

Gebe GOtt, der ein Vater ist über alles, im Himmel und auf Erden, daß alle Eltern folgen, GOtt zu Ehren, ihnen zu gutem Gewissen und endlich zu ihrer, auch ihrer Kinder Seelen Heil und Seligkeit, und helfe, daß die lieben Kinder sich auch ziehen lassen und folgen, Amen.

\chapter{Das 14. Kapitel}

\section{Vorrede}

Die Kinder sollen ihre Eltern mit Gedanken, Gebärden, Worten, Werken und Gehorsam ehren, und dies darum, daß es Gott ernstlich befohlen hat. Und wenn auch gleich kein göttlich Statut oder Gebot darauf gesetzt wäre, so bringet es doch das natürliche Recht mit sich, daß wir sollen denen Gutes und Ehre erzeigen, von denen wir alle Wohltat vielfältig empfangen haben. So ist auch gewiß, was man an den Eltern thut, das wird man auch an den Kindern zu gewarten haben. Ferner soll die Kinder, daß sie ihre Eltern ehren, bewegen, die gnädige Verheißung, welche das göttliche Statut also bald mitbringt, und sagt: darum sollst du es thun, auf daß dirs wohlgehe, und lange lebest auf Erden. Ueber das soll sie verwünschen in dem Gegenteil die schreckliche Dräuung auf den Ungehorsam, als der zeitlich und ewig mit Schand und Schaden und mit höllischem Feuer gestraft wird. Dies erinnert auch das 14. Kapitel.

\section{Simsons Heldenthat, Hochzeit und Räthsel}
\quellseite{388}
\section{Summarium}

Dies ist erstlich ein fein Exempel, wie Simson als ein gehorsam Kind, ob er wohl aus Gottes Schickung ein heidnisch Weib wollt nehmen, dennoch solches ohne seiner Eltern Vorwissen nicht thun will, V. 3.

Darnach ist eine andere Historia, wie Gott den Simson, weil er ein Fürst und Hauptmann sein sollte, der Israel aus der Philister Hand erretten sollte, mit sonderlicher Leibesstärke begnadet, daß er einen Löwen mit den Händen, ohne alle Waffen, zerreißet und erwürget, V. 6.

Zum dritten siehet man, weil Gott ihn erwecket hatte, daß er die Philister, seines Volks Israel ärgste Feinde, wohl plagen sollte, gibt ihm Gott rechtmäßige, gute und auch für der Welt genügsame Ursache, daß er sich an seinen Feinden, den Philistern, rächen kann. Einmal sie ihm und seinem Weibe dräuen, und darnach ihm das Weib nehmen, und einem andern geben, V. 20.

\section{Votum}

Der treue, barmherzige GOtt stärke alle Kinder an Leib und Seele, und regiere sie durch den heiligen Geist, daß sie ihren Eltern auch gehorsam und selige Werkzeuge sein seiner Gnaden, ihm zu Ehren und ihrem Nächsten zu Nutz und Dienst, Amen.

\chapter{Das 15. Capitel}

\section{Vorrede}

Wie über den Fürsten Simson der Geist des HErrn gerathen ist, und ihn getrieben und regieret hat: also verleihe Gott, daß wir auch Tempel und Wohnungen des heiligen Geistes sein und bleiben, der uns lehre, erinnere, stärke, leite und führe. Und wie Simson ein nützliches Werkzeug Gottes und seiner Kirchen gewesen: also helfe Gott, daß wir alle und ein jedes unter uns nach seinem Beruf und Stande ihm gehorsamlich und seiner Kirchen und derselben Gliedmaßen nützlich dienen. Wie auch Simson auf Gott mit Freuden und beständig getrauet hat und also gestorben ist: also gebe uns Gott, daß
\quellseite{389}
wir uns auch fest in wahrem Glauben an ihn halten, auf daß wir fröhlich und getrost in seinem Dienst beständig bleiben, die Trübsal uns nicht lassen abtreiben, sondern daß wir hie ritterlich ringen, durch Tod und Leben zu ihm dringen. Das verleihe uns allen der treue Gott, hochgelobet in Ewigkeit, Amen. Darauf höret ferner an das 15. Kapitel.

\section[Simson thut den Philistern durch die Füchse und]{Simson thut den Philistern durch die Füchse und den Eselkinnbacken Schaden}

a Schultern und Lenden) das ist, er schlug sie, wie sie ihm vorkamen, sie waren hohen oder niedrigen Standes.

b RamathLehi) heißt ein Hinwurf des Kinnbackens.

\section{Summarium}

Ist erstlich ein fein Exempel, wie Simson die Sachen läßt wohl an sich kommen, auf daß er desto mehr Rechts wider die Philister habe, und ihnen bezahle, was sie an den Juden verübet hatten, V. 3. Darnach siehet man, wie Vater und Tochter mit einander verbrennet werden, darum, daß der Vater die Tochter dem Simson genommen und einem andern gegeben
\quellseite{390}
hatte, V. 6. Denn solche Sünde will Gott nicht ungestraft lassen.

Zum dritten sind zwei herrliche Wunderzeichen hie, daß eine, wie Simson allein, ohne alle Waffen, tausend Mann mit einem Eselsbacken totschlägt, V. 15. Das andere, daß Gott in seiner Noth ihm einen Wasserquell läßt aus einem Backenzahn entspringen, V. 19. Solche wunderliche Historien sollen uns dahin weisen, daß wir uns Gott in allen Nöthen vertrauen und Hülfe bei ihm suchen sollen. Einmal wir sehen, daß er den Seinen wunderlich helfen kann.

\section{Votum}

Der König aller Könige, unser lieber Gott, lasse ihm alle christliche Regenten befohlen sein, und stärke und erhalte sie, daß sie auch ihm dienen und seine Kirche fördern und schützen, den armen Unterthanen zu Nutz und zu Erhaltung und Ausbreitung seines heiligen Worts, Amen.

\chapter{Das 16. Kapitel}

\section{Vorrede}

Kein Mensch ist so heilig, so verständig, so fromm, wird er sicher, und weicht Gott seine Hand von ihm, so führet und stürzt ihn der Teufel in die größte Sünde. Denn der Weg, auf dem wir allhie wandeln, ist sehr schlüpfrig, und die Anfechtung mancherlei, und der Teufel unruhig, dazu wir schwach und gar zu nachlässig. Diese menschliche Schwachheit läßt uns wohl bedenken, und derwegen wer sich läßt dünken, er stehe, der soll zuschauen, daß er nicht falle. Und wie der Herr Christus vermahnet Matth. 26, 41., so mag er wohl wachen und beten, daß er nicht in Anfechtung falle. Denn fromm sein und fromm bleiben stehet ja nicht in unserer Macht, und können es ja so bald versehen, als irgend ein anderer. O wie trefflich nöthig ist die Vermahnung und Warnung. O bester Gott, daß wir die täglich betrachten und nützlich gebrauchen, Amen. Dies soll auch erinnern das 16. Kapitel.

\section{Simsons Fall, Elend und letzte Rache}
\quellseite{391}
\section{Summarium}

Ist eine sonderliche Historia, in welcher wir sehen, wie auch heilige Leute in Sünden fallen und unrecht thun. Denn unrecht ist es, daß Simson mit den Heidinnen buhlet, V. 1. Wie wir noch sehen, daß große Herren in solchem Falle wollen mehr Freiheit haben, so doch das sechste Gebot sie eben so wohl bindet als andere Leute. Darum soll diese Historia niemand dahin deuten, daß er dergleichen auch thun möge; sondern wir sollen unsere Schwachheit an solchen Exempeln recht lernen erkennen, und desto mehr auf uns selbst Achtung haben, auch Gott mit größerem Fleiß bitten, daß er uns für Sünden und Anfechtung behüten wolle; sintemal so große Leute in grobe öffentliche Sünde gefallen, und greulich darum gestraft worden sind, wie wir weiter hören werden.

\section{Summarium}

Ist eine Historia, aus welcher wir lernen sollen, wie Gott die Hurerei an Simson gestraft hat, und seine Gnade, die er bisher gehabt, ihm entzogen, seine Augen ausgestochen, V. 21., und in einem elenden jämmerlichen Leben hat unter seinen Feinden eine zeitlang bleiben lassen. Denn wie der Psalm sagt, Gott will auch der Heiligen Sünde nicht ungestraft lassen. Nun soll man aber nicht denken (wie die Philister nicht anders wußten) daß Simson seine Stärke im Haar gehabt habe; sondern weil es Gottes Befehl war, Jud. 13, 5., daß er sollte ein Nasir, das ist, ein Verlobter Gottes sein, dem kein Schermesser auf sein Haupt sollte kommen, und doch ihm aus seinem eigenen Ansagen das Haar abgeschnitten ward, hatte er in dem Gottes Willen übertreten, und mußte also um dieser und anderer Sünde willen seiner Gnade, die er bisher gehabt, eine zeitlang beraubet sein, und mit der Stärke andern Leuten gleich werden.
\quellseite{392}
\section{Summarium}

Ist eine feine Historia, wie die Philister um ihrer Abgötterei willen gestraft werden, da sie sprachen, ihr Gott hätte ihnen ihren Feind in die Hände gegeben, V. 24. Solche Gotteslästerung müßte gestraft werden. Daß nun Simson sich so herzlich begehrt zu rächen, und Gott um Gnade der Stärke bittet, so er zuvor gehabt hatte, V. 28., solches soll man nicht dahin deuten, als wäre es erlaubt, sich selbst an seinen Feinden zu rächen. Denn Simson ist ehrlich ein Regent seines Volkes gewesen, darum hat er Recht dazu, daß er sein Volk an ihren Feinden rächte. Zum andern thut er solches nicht aus eigenem Vermögen, sondern sucht die Kraft bei Gott, und Gott läßt’s ihm auch durch Wunderwerk widerfahren.

\section{Votum}

Unser lieber GOtt und HErr regiere uns, daß wir uns für Unzucht, Abgötterei und allen andern Sünden hüten, und lindere uns die wohlverdienten Sündenstrafen, damit wir seine Gnade und Barmherzigkeit loben und preisen in alle Ewigkeit, Amen.

\chapter{Das 17. Kapitel}

\section{Vorrede}

Wenn sich die Leute für Gott, für seinem Zorn, auch für christlicher Obrigkeit Strafen nicht fürchten, so ist dem größten Haufen kein Laster und Uebelthat zu klein und geringe. Daher entstehen Abgötterei, Unzucht, Diebstahl, Mord und dergleichen Sünden. Derwegen jedermann Gott vor Augen haben soll, und die Unterthanen nicht allein gehorsam und gehorsam sein sollen ihrer Obrigkeit, sondern sollen sie auch in Ehren halten, und Gott danken, daß er Obrigkeit geordnet, die denen, so Landschaden sein, steuren, und den Frommen zu gut Gottes Wort, Friede, Zucht und Gerechtigkeit schützen und erhalten. Davon unterrichtet auch das 17. Capitel.

Micha richtet Götzendienst an.

a (Geschworen) diese Frau wird solches gelobet haben zum Bilde zu geben, das wird dem Sohn zuerst nicht gefallen haben, daß ihr darum gestohlen; darnach sich ihre Klagen lassen bewegen, und wieder gegeben, und lassen machen, was sie wollte. Da es nun wohl gerieht (wie solch Ding pflegt), ließ ers ihm auch gefallen.

b (Teraphim) deutschen wir das Heilige oder Heiligthum; wie zu unserer Zeit die hölzernen Heiligen, Todtenkränze, Todtenkleider und dergleichen heiliger Tand gehalten ist für Gottesdienst. Denn dies Wort hat unterschiedliche Bedeutung.

c Das ist, er weiset ihn; wie 2 Mos. 40, 13. stehet.
\quellseite{393}
\section{Summarium}

Ist eine Historia, wie sich im Stamm Ephraim eine neue Abgötterei so liederlich hat angefangen, und ein Levit sich dazu gefunden, der Armuth halben, um des Bauchs willen, zu solcher Abgötterei sich hat brauchen lassen, V. 6. Und meldet der Text sonderlich, daß solches zu der Zeit geschehen sei, da kein König noch Regiment im Lande gewesen sei, V. 7. Denn weltliche Obrigkeit soll fürnehmlich darauf sehen und halten, daß rechte Lehre und rechter Gottesdienst erhalten, und aller Abgötterei gewehret werde. Wo nun das Regiment gefallen ist, da muß die Kirche auch mit erhalten, und sinket der Satan Mittel und Wege, wie er Irrthum und Abgötterei anrichte. Derohalben jedermann desto fleißiger für fromme Regenten und ihr langes Leben bitten solle.

\section{Votum}

Der treue GOtt erhalte unsere liebe Obrigkeit in langem Leben und beständiger Gesundheit, und verleihe, daß ihre Regierung ferner glücklich und nützlich sei, auf daß sein Lob, Güte und Wohlthaten unter uns allzeit gerühmet und gepreiset werden, bis an den jüngsten Tag, Amen.

\chapter{Das 18. Kapitel}

\section{Vorrede}

Sirach im 15. Kapitel V. 13 spricht: Der HErr hasset alle Abgötterei, und wer ihn fürchtet, der scheuet sich dafür. Das Buch der Weisheit im 14. Kapitel V. 27. lehret: Den sündlichen Götzen dienen, ist alles Bösen Anfang, Ursache und Ende. Und St. Paulus sagt 1 Corinth. 6, 9.: Die Abgöttischen werden das Reich Gottes nicht ererben. Diese göttlichen Donnerschläge sollten uns ja genug sein von Abgötterei abzustehen. Aber da ein Lämmlein für dem Wolf erschrickt, und da ein Täublein für dem Geier sich fürchtet, so laufen die Menschen, ja die ausbündige Christen sein wollen, der Abgötterei zu und helfen die bestätigen. Derwegen wende sich ein jeder vom Unrechte und hasse mit Ernst die Abgötterei, und dagegen haltet fleißig an rechter Lehre und rechtem Gottesdienst, so wird es Gott wieder mit euch halten, zu zeitlicher und ewiger Wohlfahrt. Darauf höret auch davon an das 18. Kapitel.

\section{Die Daniter gewinnen Lais, und richten Abgötterei an}

a Sie hatten weder Fürste noch Herren.
\quellseite{394}
\section{Summarium}

Ist eine Historia, wie die Abgötterei aus dem Hause Micha unter den ganzen Stamm der Kinder Dan kommen sei, V. 13., auf daß wir an solchen Exempeln lernen, daß wir bei Leib der Abgötterei und falscher Lehre nicht lang zusehen, sondern je eher je besser dazu thun und ausrotten; denn beides, die falsche Lehre und Abgötterei, hat die Art, wie St. Paulus sagt, daß es kreucht, und frißt wie der Krebs, und immer viel Leute und Anhangs zu sich zeucht. Item, es durchdringet wie der Sauerteig, darum ists Noth, daß man alsobald wehre; und gleich wie das Feuer nicht lasse Luft gewinnen, sonst ists hernach ungewehret, wenns überhand nimmt.

\section{Votum}

Wir danken Gott herzlich, der uns von der antichristlichen Abgötterei erlöset hat; der verleihe, daß wir auch alle Abgötterei allzeit ernstlich meiden, damit wir uns und gemeinem Vaterlande nicht Gottes gerechte Strafen zuziehen, Amen.
\quellseite{395}
\chapter{Das 19. Kapitel}

\section{Vorrede}

Gott hat nächst dem Unglauben keine Sünde und Laster so hart und greulich gestraft, und weniger ungestraft wollen lassen hingehen, als Hurerei und Ehebruch. Was ist die Ursache, daß die Sündflut über die Welt und das Schwefel, Feuer und Pech über Sodom kam und regnet? Keine andere, denn die Sünde wider den heiligen Ehestand, nämlich Unzucht und Ehebrecherei. Warum wurden 24,000 Israeliten in der Wüste, alle Midianitische Weiber und hernach 25,000 Benjaminiter erwürget? Um der greulichen Unzucht willen. Diese und dergleichen erschreckliche Strafen Gottes sollen Junge und Alte von aller Unreinigkeit in und außerhalb dem Ehestande abschrecken. Davon unterrichtet auch das 19. Kapitel.

\section[Die zu Gibea begehen eine schändliche und höchst]{Die zu Gibea begehen eine schändliche und höchst sträfliche That}
\quellseite{396}
\section{Summarium}

Ist eine schreckliche und greuliche Historia von den Gibeonitern, wie sie den Leviten und sein Weib nicht beherberget, sondern ihr alle Schand und Schmach angelegt haben, V. 22. Solche Exempel sollen uns zur Furcht Gottes und fleißigem Gebet vermahnen, daß uns Gott für Sünden und Schanden behüten und in Anfechtung nicht wolle sinken lassen. Denn wo Gott nicht hält und behütet, ist's sehr bald geschehen, daß man in Schand und Laster fällt; wie wir an Simson, David und andern Heiligen sehen. Darum Christus der HErr im Vater unser uns beten gelehret: Führe uns nicht in Versuchung.

\section{Votum}

Verleihe der treue, liebe GOtt, daß wir folgen; dem Teufel und unsers Fleisches bösen Lüsten, auch der Welt Anreizung Widerstand thun und in Heiligkeit und Reinigkeit allzeit leben und auch sterben, Amen.

\chapter{Das 20. Kapitel}

\section{Vorrede}

Gott läßt die Unzucht nicht ungestraft, auch an denen, die Buße thun und sich nach begangener Unzucht bekehren. Warum lassen sich denn heut so viele unter Jungen und Alten, in und außerhalb dem Ehestande, den unsaubern Geist so verblenden und verführen? Demgegen wollen wir beten, Gott der Allmächtige, der ein reines, keusches Wesen ist, und auch Reinigkeit und Keuschheit von uns haben will, wolle durch sein Wort und durch seinen heiligen Geist unsere Herzen reinigen von allen unreinen Worten, Werken und Gedanken, und daß er uns gebe, daß wir ihm mit reinem Herzen dienen, ihn anrufen und ihm danken, hie zeitlich und dort ewiglich. Denn selig sind, die da reines Herzens sind, denn sie werden Gott schauen, sagt Christus Matth. 5, 8. Das soll auch erinnern das 20. Kapitel.
\quellseite{397}
a Non petunt victoriam, sed viribus praesumunt, et justitia sua confidunt, ac animantur malitia filiorum Benjamin. D. i. sie bitten nicht um Sieg, sondern verlassen sich auf ihre Kraft und vertrauen auf ihre Gerechtigkeit, und werden muthig durch die Bosheit der Kinder Benjamin.

b Weil die Kinder Israel auch wohl diesmal nicht um Sieg den Herrn bitten, so verheißet er ihnen auch den Sieg nicht.

\section{Summarium}

Dies ist eine Historia, in welcher man siehet, wie oft der schuldige und ungerechte Haufe zu einer bösen Sache Glück hat, wider den andern Theil, der doch der Sachen Recht hat, aber sonst bußwürdig ist. Also hatten die Benjamiter gröblich Unrecht, daß sie die greuliche Sünde der Gibeoniten nicht allein nicht strafen, sondern die Uebelthäter auch vertheidigen lassen. Wiederum die andern elf Stämme thaten recht daran, daß sie die greuliche Sünde nicht wollen ungestraft lassen; und dennoch liegt Benjamin zweimal ob, und die andern elf Stämme verlieren vierzig tausend Mann in einer solchen
\quellseite{398}
guten Sache, 20, 24. Warum gehet es doch so ungleich zu? Sind doch die Gibeaen und Benjaminiten sträflich. Die elf Stämme aber haben guten Fug und Recht, und liegen dennoch unten. Woher kömmt doch solches? Antwort: Obgleich die Sache an sich selber nicht böse war, so sind doch die elf Stämme in andern Fällen sträflich gewesen, nämlich aufs erste in dem, daß sie den Sieg von Gott nicht baten, sondern sich verließen auf beides, ihre Macht und gute Sache, und daß ihr Gegenteil unrecht hatte. Darnach auch, daß sie, wie oben stehet, öffentlich Abgötterei lang her getrieben und geduldet hatten, daher ist solches Urtheil über sie gangen.

c Ich, spricht Gott, will es thun. Ziehet habt ihrs wollen thun; aber es heißet, ich, nicht ihr.

d Der Herr schlug sie, nicht Israel, denn es heißt: gloria in excelsis Deo.
\quellseite{399}
\section{Summarium}

Sie ist eine Historia, in welcher man stehet, obschon Gott bisweilen bösen Buben etwas lang zusiehet, daß er dennoch endlich mit der Strafe kommet und so viel desto härter zuschlägt. Sonderlich aber soll man nicht vergessen, daß die greuliche Unzucht, welche die Gibeoniter begangen hatten, und die Benjamiten vertheidigten, so jämmerlich und hart gestrafet wird, daß alle Städte in Benjamin verheeret, und der ganze Stamm fast gar ausgerottet ward. Denn Hurerei und Unzucht ist auch der Sünden eine, welche Gott mit Verwüstung von Land und Leuten strafen will, wo weltliche Obrigkeit nicht strafet; wie man hie und an andern Orten siehet. Darum jedermann sich mit höchstem Fleiß dafür hüten, und weltliche Obrigkeit mit harter Strafe solchem Aergerniß wehren soll.

\section{Votum}

Solches helfe der treue Gott uns und aller Obrigkeit, daß es durch seine Gnade und Kraft des heiligen Geistes geschehen möge, zur Abwendung aller Strafen, zeitlich und ewiglich, Amen.

\chapter{Das 21. Kapitel}

\section{Vorrede}

Wie eine Biene, wenn sie aus Zorn den Menschen sticht, darüber ihren Stachel verlieret, keinen Honig mehr machen kann, sondern ein steter Knecht sein, und Wasser zutragen muß: also wirket des Menschen Zorn nichts Gutes. Da vergehet alle Freude und Lust etwas Gutes zu thun, und wird der zornige Mensch ein Knecht des Teufels und der Sünden, schadet seinem Nächsten so sehr nicht als ihm selbst. Und wie ein Wasser sich selbst versiegen und endlich gar verraucht, also verlieret ein Zorniger auch die Gnade des heiligen Geistes. So ist der Zorn ihnen eine Thür und Schrotleitern zu allerlei Untugend und Schandlustern. Derwegen ein jeder Christ für ungerechtem Zorn sich fleißig hüten soll. Davon höret auch an das 21. Kapitel.

\section{Wie der Stamm Benjamin wieder erlaubet werden}
\quellseite{400}
\section{Summarium}

Dies ist eine Historia, wie die elf Stämme die übrigen vom Stamm Benjamin wieder zu Gnaden annehmen, und sie in das Ihrige einsetzen, ihnen helfen und rathen. Denn es soll mit der Strafe und dem Zorn ein Maß haben, daß man ihm nicht zu viel thue.

\part{Das Buch Ruth}

\chapter{Das 1. Kapitel}

\section{Vorrede}

Dies Buch Ruth, neben dem, daß es eine feine, liebliche und sonderlich zum Hausregiment nützliche Historia ist, dienet auch dazu, daß wir lernen, wie Gott sich zur selbigen Zeit auch der Heiden hat angenommen, als die hernach zu seinem Reich durchs Evangelium auch kommen, und des Samens Abrahams genießen sollten. Und ist sonderlich hie wohl zu merken, daß diese Ruth, eine Heidin aus der Moabiter Land, unter den Voraltern des HErrn Christi gezählet wird, wie auch Matth. 1. v. 5. zu sehen ist.

Im ersten Kapitel soll man lernen, wie eine Schnur sich gegen ihre Schwieger freundlich, demütig und dienstlich halten soll. Denn solches gefällt Gott wohl, will es auch reichlich vergelten, wie er an der frommen Ruth bewiesen hat, und wir darnach weiter hören werden.

\section{Ruth reiset mit Naemi nach Bethlehem}
\quellseite{401}
a) Naemi) heißt meine Lust.

b) Mara) heißt bitter oder betrübt.

\section{Summarium}

Bei diesem Kapitel sollen wir merken drei Stücke. Das erste ist von der Theurung im jüdischen Lande, welche Gott über Israel hat ergehen lassen von wegen ihrer mannigfältigen Sünden; denn ein jeder gelebet hat, wie es ihm gelüstet, weil sie etliche Zeit keinen Fürsten und Regenten gehabt hatten. Also straft Gott noch die Sünder der Welt, wie mit andern Plagen, also mit Theurung, wie denn Gott zum oftermal gedräuet hat, er wolle über die Sünder seinen Himmel zuschließen. Das andere Stück ist, daß Elimelech mit den Seinen kehret, daß man bei Abgöttischen und Gottlosen also wohnen soll, daß solche Beiwohnung uns gleichwohl an unserm Glauben und Religion nicht schädlich sei. Aber zu unsern Zeiten werden viele Leute gefunden, welche die Religion mit Veränderung ihrer Wohnung auch verändern. Daher etliche das Sprüchwort haben: Darum ich wohne, darnach glaube ich. Das dritte Stück ist, Naemi heißt ihre beiden Schnuren wiederum heimkehren, damit sie bei den Juden, als bei einem fremden Volke, nicht in Gefahr des Lebens und der Ehren kommen. Arpa thut's, Ruth aber bleibt bei ihrer Schwieger. Das ist ein schön Exempel, daß die Schwiegerfrauen und Schnuren mit einander treulich haus
\quellseite{402}
beln, und einander herzlich lieben, folgen und dienen sollen. Und daß Naemi ihre Schwuren segnet, und wünschet ihnen Gutes, das denn die Ruth hernach bekommen hat, lehret, daß der Eltern Segen über ihre ungehorsame Kinder nicht vergeblich ist. Wie denn auch Sirach saget Cap. 3. V. 10.: Der Segen der Eltern erbauet die Häuser der Kinder.

\section{Votum}

Der treue GOtt lindere uns die wohlverdienten Strafen, und erhalte uns bei der christlichen Religion und in herzlicher Liebe gegen Ihm, gegen den Unsern und gegen unsern Nächsten, auf daß wir einander gern und treulich dienen, Amen.

\chapter{Das 2. Capitel}

\section{Vorrede}

Wer etwa ein Fremdling ist, und wollte gern Gnade und Gunst bei fremden Leuten haben, der besiehe sich, daß er durch den Glauben einen gnädigen Gott habe, und daß er christlich lebe; so wirds Gott also schicken und machen, daß er mehr Liebe und Huld bei Fremden finden wird, denn bei seinen Freunden und Verwandten. Also wer dem Nächsten Gutes thut, ob er wohl das Guts wünschen und Danksagen der Leute dafür nicht höret, so schreibet doch das Werk an ihm selbst, da es aus dem Glauben geschieht, zu Gott. Denn der Glaube ist gleich als ein Wagen, dadurch alle unsere Werke für Gottes Angesicht geführet werden, der sie auch belohnet, wie Sirach sagt im 17. Capitel V. 18.: Gott behält die Wohltat der Menschen wie einen Siegelring, und die guten Werke wie einen Augapfel. Dies sind Exempel Ruth und Boas. Davon das 2. Kapitel.

\section{Ruth lieset auf Boas Felde Aehren auf}

a Das ist, sie ist nicht der Wehen eine, die daheim auf dem Polster sitzen und faullenzen.
\quellseite{403}
\section{Summarium}

Dies ist ein fein Exempel, erstlich an Ruth, wie sie Gott der Treue, so sie an ihrer Schwieger gethan, genießen läßt, V. 3. Darum muß folgen, wo Schnur und Schwieger, wie gemeiniglich geschieht, unter einander unfreundlich und undienstlich sind, daß Gott übel gefalle, und darum armselige, elende Leute müssen sein.

Zum andern ist ein sehr fein Exempel an Boas, der seine arme Freundin oder Schwägerin nicht verachtet, sondern ist ihr freundlich, und thut ihr alles Gutes, V. 8. Denn Gott gibt oft manchem viel Segens und Guts, allein, daß er seinen armen Freunden damit helfen, und sie erfreuen und trösten soll. Wo mans aber nicht thut, und sich armer Freunde nicht annimmt, und ihnen nicht gern hilft, solche Untugend wird Gott zu seiner Zeit wohl finden und strafen.

\section{Votum}

Verleihe GOtt, daß Freundschaft also hie gegen einander lebe, daß sie auch in ewiger Freundschaft beisammen leben; und helfe allen, tröste und fördere auch alle, um Christi JEsu willen, die armen und verlassenen Christen mit Rath und That aushelfen, Amen.

\chapter{Das 3. Kapitel}

\section{Vorrede}

Um den Ehestand ist es ein langer Kauf, darinnen erneuern und hinter sich sehen nicht gilt, noch hilft. Darum mag man wohl langsam, fürsichtig und weißlich damit umgehen, und den Ehestand in aller Furcht Gottes, und mit wohlbedachtem Muth fürnehmen. Denn er ist nicht ein Pferdekauf, sondern was einmal geschlossen wird, dabei muß man folgend beruhen. Dies sollen wohl betrachten beides Junggesellen und Jungfrauen. Denn es heißt: Ehe wiegs, denn wags. Und ein tugendsam Eheweib und frommer Mann kömmt und wird gegeben von dem HErrn, der ist Gottes Gabe und Segen, wie Salomon in seinen Sprüchen im 19. Kapitel V. 14. lehret. Und derwegen mögen junge Leute wohl beten, Augen und Ohren aufthun, sehen, hören und bedenken, was sie thun, damit sie es nicht hernach gereue. Ferner höret darauf an das 3. Kapitel.

\section{Ruth bekommt wegen künftiger Heirath guten Bescheid}
\quellseite{404}
\section{Summarium}

Ist eine feine Historia, wie die arme Ruth aus Befehl ihrer Schwieger ihren nächsten Freund Boas um die Ehe in aller Zucht und Demuth anspricht, V. 6. Und Boas fürnehmlich bedenket, daß sie sich züchtig und wohl gehalten, ein gutes Geschrei bei jedermann erhalten, gottesfürchtig und fromm gewesen sei, V. 12. Darum sollen hie beide Jungfrauen und junge Wittwen lernen, daß sie gottesfürchtig, züchtig und fromm sein, so werden sie es genießen. Desgleichen die jungen Gesellen hie sollen auch ein Exempel nehmen an Boas, daß sie, wenn sie freien wollen, mehr auf Gottesfurcht, Zucht und Frömmigkeit, denn auf Geld und Gut sehen sollen. Weil aber solches selten geschieht und wenig vermaßen in den Ehestand sich schicken, siehet man, daß desto mehr Unglücks und Unraths unter den Leuten ist.

\section{Votum}

Der Stifter des heiligen Ehstandes, unser lieber GOtt, helfe, daß junge Leute in diesen Orden in seiner Furcht, mit vorhergehendem herzlichen Gebete und mit gutem Bedacht treten, und daß sie darinnen christlich leben, und Segen an Leib, Seele und Nahrung haben und behalten, Amen.

\chapter{Das 4. Kapitel}

\section{Vorrede}

Jetzt sind die Leute gemeiniglich so geschickt, daß sie heimlich, ohne andrer Leute Beisein, auch mit Geldbuhr zusagen, das verleugnen sie hernach öffentlich. Also Geizhälse suchen allein ihrem Frommen und Nutzen nach, und lassen andern Leuten Schaden. Ferner, daß Gott Kinder bescheret, dieser Segen wird jetzt schier vor einen Fluch gehalten, welches daher kömmt, daß ihrer viel die Menge der Kinder nicht ansehen nach Gottes Wort durch den Glauben, sondern nach ihren Gedanken auf die Beschwerungen, die sich durch die Kinder und Kinderzucht zutragen. Anders lehret davon das 4. und letzte Kapitel.

\section{Die Hyrath Boas mit Ruth wird vollzogen und gesegnet}

a Wo du hin willst, denn das gebräuchliche Ploni almoni ist ein Name eines ungewissen Orts.
\quellseite{405}
b) Id est, Det Deus, ut cum illa magnificeris: ut certe factum est. Nam peperit Obed, avum Davidis. Das ist, (Gebe Gott, daß du mit ihr groß werdest; wie auch wirklich geschehen ist, denn sie hat den Obed geboren, den Großvater Davids.

\section{Summarium}

Dies ist eine Historia, wie Boas die arme Wittwe Ruth, seine Freundin, zur Ehe nimmt, V. 10. Und Gott den frommen und gottesfürchtigen Leuten einen Sohn gibt, welcher Davids Großvater gewesen ist; auf daß wir lernen, wie Gott Zucht und Frömmigkeit lieb hat und segnet, dagegen unzüchtige, furchtlose Leute weder Glück noch Heil haben, oder es zu ihrem großen Schaden und Nachtheil haben.

\section{Votum}

Gott, der HErr, verleihe, daß alle Eheleute auch in Gottes Furcht und in Zucht leben, und gebe ihnen seinen reichen Segen an Kindern, Gut und Nahrung, doch nach seinem Willen, und wie es den Eheleuten gut und selig sein soll, Amen.
